\documentclass[11pt,a4paper]{article}
\usepackage[T1]{fontenc}
\usepackage[utf8]{inputenc}
\usepackage{amsmath,amsthm,mathtools}
\usepackage{newpxtext,newpxmath}
\usepackage[a4paper,margin=25mm,headheight=15pt]{geometry}
\usepackage{microtype}
\usepackage{xcolor,booktabs,array,longtable}
\usepackage{enumitem}
\usepackage{fancyhdr}
\usepackage{listings}
\usepackage[most]{tcolorbox}
\usepackage{hyperref}
\definecolor{ink}{HTML}{183547}
\definecolor{accent}{HTML}{126B73}
\definecolor{pale}{HTML}{EEF5F6}
\hypersetup{colorlinks=true,linkcolor=accent,citecolor=accent,urlcolor=accent,
  pdftitle={Guarded periodic blocks and exact maximal order types of words below omega squared},
  pdfauthor={ChatGPT, research report prepared for Vladimir Reshetnikov},
  pdfsubject={A proposed solution to Altman's finite-alphabet sharpness question}}
\pagestyle{fancy}
\fancyhf{}
\fancyhead[L]{\small\textsc{Guarded periodic blocks}}
\fancyhead[R]{\small\textsc{Exact order types below $\omega^2$}}
\fancyfoot[C]{\thepage}
\renewcommand{\headrulewidth}{0.3pt}
\setlength{\parskip}{3pt}
\setlength{\parindent}{1.2em}
\setlength{\emergencystretch}{2em}
\setlist{itemsep=2pt,topsep=4pt}
\numberwithin{equation}{section}
\newtheorem{theorem}{Theorem}[section]
\newtheorem{lemma}[theorem]{Lemma}
\newtheorem{proposition}[theorem]{Proposition}
\newtheorem{corollary}[theorem]{Corollary}
\theoremstyle{definition}
\newtheorem{definition}[theorem]{Definition}
\newtheorem{example}[theorem]{Example}
\theoremstyle{remark}
\newtheorem{remark}[theorem]{Remark}
\newcommand{\mot}{\mathop{o}}
\newcommand{\W}{\mathcal W}
\newcommand{\J}{\mathcal J}
\newcommand{\F}{\mathcal F}
\newcommand{\G}{\mathcal G}
\newcommand{\eps}{\varepsilon}
\newcommand{\len}{\mathop{\rm len}}
\newcommand{\down}{\mathord{\downarrow}}
\newcommand{\emb}{\preccurlyeq}
\newcommand{\eqv}{\simeq}
\newcommand{\N}{\mathbb N}
\newcommand{\Hn}[1]{H_{#1}}
\newcommand{\Wordspace}[1]{s_{\omega^2}^{F}(#1)}
\newcommand{\E}{\mathcal E}
\newcommand{\Max}{\operatorname{Max}}
\newcommand{\rec}{\operatorname{Rec}}
\newtcolorbox{resultbox}{enhanced,breakable,colback=pale,colframe=accent,
  boxrule=0.6pt,arc=1mm,left=3mm,right=3mm,top=2mm,bottom=2mm}
\lstset{basicstyle=\ttfamily\small,breaklines=true,columns=fullflexible,
  frame=single,rulecolor=\color{accent},backgroundcolor=\color{pale},
  showstringspaces=false,tabsize=4}
\IfFileExists{data/test_summary.tex}{\newcommand{\TestAssertions}{3,240,040}
\newcommand{\ProductChecks}{1,625,127}
\newcommand{\MarkerStages}{355}
\newcommand{\DeciderChecks}{774,149}
\newcommand{\BinaryPairs}{726,242}
\newcommand{\BinaryStages}{14}
}{
\newcommand{\TestAssertions}{3,240,040}
\newcommand{\ProductChecks}{1,625,127}
\newcommand{\MarkerStages}{355}
\newcommand{\DeciderChecks}{774,149}
\newcommand{\BinaryPairs}{726,242}
\newcommand{\BinaryStages}{14}}

\title{\vspace{-10mm}\color{ink}\LARGE\bfseries
Guarded periodic blocks and exact maximal order types\\[2mm]
of words below $\omega^2$\\[3mm]
\large A proposed solution to Altman's finite-alphabet\\sharpness question}
\author{Research report prepared with ChatGPT\\
\small For Vladimir Reshetnikov}
\date{September 19, 2026}

\begin{document}
\maketitle
\thispagestyle{empty}
\begin{abstract}
For a finite poset $P$, consider all words of ordinal length less than
$\omega^2$, ordered by strictly increasing subsequence embeddings that may
increase labels in $P$. Write $n=|P|$ and let $j(P)$ count all downsets of
$P$, including the empty one. This article gives a complete proposed proof
that
\[
  \mot\bigl(\Wordspace P\bigr)=
  \begin{cases}
    1,&n=0,\\
    \omega^2,&n=1,\\
    \omega^{\omega^{\,n+j(P)-2}},&n\ge2.
  \end{cases}
\]
In particular, the two-letter antichain has maximal order type
$\omega^{\omega^4}$, answering the explicit sharpness question in
Example~3.15 of Altman's March 2026 revision. The two-letter chain has type
$\omega^{\omega^3}$. The main construction embeds every finite Cartesian
power of an earlier stage into the next stage by synchronizing periodic
$\omega$-blocks. Proper downsets admit an excluded-letter guard; the universal
block instead requires a limit-length padding argument. A more general
classification for selected families of periodic tails isolates the exact
exception in which naive block counting fails. Two routes from the product
embeddings to the ordinal lower bound are given: the classical natural-product
theorem, and an elementary hand-built linear extension of a Cartesian power
that makes the main theorem independent of that import. Supplementary
exact-length strata are computed, and they show explicitly why the answer
cannot be assembled length by length. An exact symbolic embedding algorithm
with canonical normal forms, negative controls for three invalid shortcuts,
and a reproducible test suite accompany the proofs.
\end{abstract}

\begin{resultbox}
\textbf{Status and scope.} This is a new, unrefereed research argument, not a
report of an independently certified result, and it has not been checked in a
proof assistant. All proof steps beyond the explicitly stated classical
maximal-order-type theorems are supplied below. The specific unresolved
question was checked in the cited source; the claim of a solution rests on the
argument given here and not on a claim of exhaustive literature coverage.
The computer checks test symbolic embeddings, not infinite ordinal ranks.
The theorem concerns \emph{strictly increasing} embeddings, finite alphabet
posets, and the strict length cutoff $\omega^2$; it does not settle arbitrary
alphabets or higher cutoffs.

\medskip
\textbf{Provenance.} This report merges two independently produced research
reports on the same question, \texttt{order-types-below-omega-squared} and
\texttt{guarded-periodic-blocks}. Both reached the same formula by the same
core mechanism: a finite generating alphabet for the upper bound, cofinal
synchronization of periodic separators, an excluded-letter guard for proper
downsets, and a limit-length guard for the universal one. That shared core is
proved exactly once below. The two reports diverged at one step, the passage
from the Cartesian-power embeddings to the ordinal lower bound, and both
arguments are kept: the classical natural-product route in
Corollaries~\ref{cor:properlevel} and~\ref{cor:fulllevel}, and the elementary
route of Section~\ref{sec:altroute}. Two different justifications of the
monotone-surjection principle are likewise both recorded in
Section~\ref{sec:mot}. Material contributed by the second report, and absent
from the first, includes the exact-length strata of
Section~\ref{sec:exact}, the canonical normal forms of
Section~\ref{sec:canonical}, and the worked negative controls of
Section~\ref{sec:negative}.
\end{resultbox}

\vfill
\noindent\small\textbf{Keywords.} Well-quasi-order; maximal order type;
transfinite word; recurrent downset; guarded periodic block; Higman ordering;
finite downset lattice; ordinal arithmetic; natural product.
\newpage
\begingroup
\small
\setlength{\parskip}{0pt}
\tableofcontents
\endgroup
\newpage

\section{The selected question and the result}
\label{sec:intro}

Altman studies finite-image words of bounded transfinite length in
\cite{Altman2026}. Example~3.15, on printed page~10 of version~2, gives the
upper bound $\omega^{\omega^4}$ for a two-element antichain and explicitly
leaves its sharpness unverified. The same example improves the chain's
upper bound to $\omega^{\omega^3}$. The arXiv record checked for this note
lists version~2, dated March~10, 2026, as the latest revision. Targeted
searches did not locate a subsequent resolution; this is a limited
literature check, not a guarantee of priority. No claim is made that the
present literature check excludes every earlier formulation of the general
finite-poset formula.

\paragraph{Related work.}
The classical maximal-linear-extension theory is due to de Jongh and
Parikh~\cite{deJonghParikh1977}, and the finite-word calculation is part of the
de Jongh--Parikh--Schmidt theory \cite{deJonghParikh1977,Schmidt2020}; the
underlying well-quasi-order assertion is Higman's lemma \cite{Higman1952}.
Recent work of Chopra and Pakhomov \cite{ChopraPakhomov2026} studies the limit
cutoff $\omega^\omega$ through finite trees. The present article addresses the
finite cutoff $\omega^2$ and the explicitly identified two-letter example, so
the two lines of work are complementary rather than overlapping.

This is an attractive problem to attack because the alphabet is tiny, all
$\omega$-tails have finite descriptions up to embedding, and the candidate
ordinal is explicit. The difficult part is not finding more tail types. It
is proving that the types already available contribute their full ordinal
complexity after concatenation.

Let $P$ be a finite poset. A \emph{downset} is a subset $D\subseteq P$ such
that $x\le_P y\in D$ implies $x\in D$. Put
\[
 \J(P)=\{D\subseteq P:D\text{ is a downset}\},\qquad j(P)=|\J(P)|.
\]
The empty set is counted in $j(P)$ throughout.

\begin{theorem}[Finite-alphabet formula]\label{thm:main}
Let $P$ be a finite poset with $n=|P|$ elements. With strictly increasing,
label-respecting subsequence embedding and mutual embeddability quotiented
out,
\begin{equation}\label{eq:main}
 \boxed{\mot\bigl(\Wordspace P\bigr)=
 \begin{cases}
  1,&n=0,\\
  \omega^2,&n=1,\\
  \omega^{\omega^{\,n+j(P)-2}},&n\ge2.
 \end{cases}}
\end{equation}
\end{theorem}

The two small alphabet sizes are part of the theorem, not removable
technicalities. In particular, substituting $n=1$ into the third line would
give an incorrect answer; Section~\ref{sec:singleton} explains why the defect
is structural.

Here $s_{\omega^2}^{F}(P)$ denotes words of length \emph{strictly less than}
$\omega^2$. Since $P$ is finite, the finite-image superscript is automatic.
The theorem is proved in Section~\ref{sec:classification} after two product
embedding lemmas.

\begin{corollary}[The selected binary cases]\label{cor:binary}
For a two-element antichain $A_2$ and a two-element chain $C_2$,
\[
 \mot\bigl(\Wordspace {A_2}\bigr)=\omega^{\omega^4},
 \qquad
 \mot\bigl(\Wordspace {C_2}\bigr)=\omega^{\omega^3}.
\]
\end{corollary}
\begin{proof}
There are four downsets in $A_2$ and three in $C_2$. Substitute $n=2$ in
\eqref{eq:main}.
\end{proof}

The proof establishes more than these two values. We may choose an arbitrary
family of nonempty downsets and allow only the corresponding periodic
$\omega$-blocks. Unless the universal block is the sole allowed infinite
block, the number of allowed block types determines the maximal order type
exactly. The exceptional case is essential, not a technical inconvenience.

\subsection{What the proof adds}
The upper bound follows from a finite generating alphabet. The lower bound
requires an additional construction. A new periodic block serves as a
separator only when an embedding cannot hide it inside earlier blocks.
Introducing downsets in inclusion order solves that problem. A finite guard
then prevents a component from leaking into the next separator. At the
universal separator no excluded letter exists, so a different device is
needed: attach a cancellable tail of limit length. An infinite tail cannot
fit into a bounded, hence finite, initial segment of an $\omega$-block.

Together these observations give actual order embeddings
\[
  X^{r+1}\hookrightarrow Y\qquad(r<\omega)
\]
from each stage into the next. Classical natural-product arithmetic turns
these embeddings into the sharp lower bound. Section~\ref{sec:altroute} gives
a second, elementary passage from the same embeddings to the same bound, using
only an explicitly constructed linear extension of a Cartesian power; the
main theorem therefore does not depend on importing the de Jongh--Parikh
product formula.

\paragraph{What is claimed as the contribution.}
The normal-form encoding and the finite Higman maximal-order-type formula are
established ingredients. The central contribution is the guarded product
construction of Sections~\ref{sec:proper} and~\ref{sec:universal}, which makes
the generator-count upper bound sharp. The support-family classification and
the finite-poset formula follow from that construction. The exact-length
calculations of Section~\ref{sec:exact} are supplementary consequences, not
the foundation of the main lower bound.

\section{Definitions and classical inputs}
\label{sec:prelim}

\subsection{Ordinal words and embeddability}
A \emph{word over $P$} is a function $u:\alpha\to P$, where $\alpha$ is an
ordinal; its length is $\len(u)=\alpha$. If $u:\alpha\to P$ and
$v:\beta\to P$, write $u\emb v$ when there is a strictly increasing map
$f:\alpha\to\beta$ satisfying
\[
 u(\xi)\le_P v(f(\xi))\qquad(\xi<\alpha).
\]
No continuity or cofinality condition is imposed on $f$. Write $u\eqv v$
when $u\emb v$ and $v\emb u$. Concatenation is ordinal concatenation. The
empty word $\eps$ has length zero. For lengths below $\omega^2$ there are
unique $k,m<\omega$ such that
\[
 \len(u)=\omega k+m.
\]

Embeddability is in general only a quasi-order. Distinct words, such as
$a^\omega$ and $aa^\omega$, may be equivalent. All maximal order types in
this article refer to the quotient by $\eqv$.

Strictness of the index map is indispensable. Under a merely nondecreasing
index map an infinite constant source could be collapsed onto a single target
position, and every statement below would describe a different problem. The
same remark explains why finite material preceding an infinite block can
disappear from the \emph{length} expression, by $k+\omega=\omega$, while its
labels may still matter for embeddability.

\begin{lemma}[Concatenation is monotone]\label{lem:concat}
If $u_i\emb v_i$ for $0\le i<r$, where $r$ is finite, then
$u_0\cdots u_{r-1}\emb v_0\cdots v_{r-1}$.
\end{lemma}
\begin{proof}
Translate the embedding on each component into the corresponding interval
of the target concatenation. The intervals occur in the same order, so the
union of these translated maps is strictly increasing and respects labels.
\end{proof}

A finite concatenation of words of length below $\omega^2$ still has length
below $\omega^2$. In particular, all maps used later stay within the stated
cutoff, even though the number of $\omega$-blocks in their arguments is
unbounded across different arguments.

\subsection{Maximal order types}
\label{sec:mot}
A quasi-order is a \emph{well-quasi-order} if every infinite sequence
$x_0,x_1,\ldots$ contains $i<j$ with $x_i\le x_j$. The maximal order type
$\mot(X)$ is the largest ordinal type of a linear extension of the quotient
poset. The existence of this maximum, its characterization by the tree of
bad finite sequences, and the product theorem below are classical
\cite{deJonghParikh1977}. We use the following precise inputs.

\begin{theorem}[Classical maximal-order-type calculus]\label{thm:classical}
For well-quasi-orders $X,Y$:
\begin{enumerate}[label=\textnormal{(\roman*)}]
\item An order embedding $X\hookrightarrow Y$ implies $\mot(X)\le\mot(Y)$.
\item An order-preserving surjection from $Y$ onto $X$, including surjectivity
up to quasi-order equivalence, implies $\mot(X)\le\mot(Y)$.
\item For the componentwise Cartesian product,
$\mot(X\times Y)=\mot(X)\otimes\mot(Y)$, where $\otimes$ is natural product.
\item For a disjoint union, with no comparisons across the two parts,
$\mot(X\amalg Y)=\mot(X)\oplus\mot(Y)$, where $\oplus$ is natural sum.
\item If $P$ is a nonempty finite poset of cardinality $n$, its Higman word
ordering satisfies
\begin{equation}\label{eq:finiteword}
 \mot(P^*)=H_n,\qquad H_n:=\omega^{\omega^{n-1}}.
\end{equation}
\end{enumerate}
\end{theorem}

The finite-word formula is part of the de Jongh--Parikh--Schmidt theory
\cite{deJonghParikh1977,Schmidt2020}; its underlying well-quasi-order
assertion is Higman's lemma \cite{Higman1952}. These are imported theorems;
the present article does not claim new proofs of them. All later lower bounds
are reduced to these inputs by explicit embeddings. Assertions~(iii),
(iv) and~(v) are used as follows: (v) everywhere, (iii) in the main line of
Corollaries~\ref{cor:properlevel} and~\ref{cor:fulllevel} and, together
with~(iv), in the exact-length calculations of Section~\ref{sec:exact}. The
alternative route of Section~\ref{sec:altroute} shows that (iii) can be
dropped from the main theorem, though not from Section~\ref{sec:exact}.

Assertion~(ii) admits two elementary explanations, and both are recorded
because they illuminate different sides of the statement.

\emph{First argument, on bad sequences.}
Choose a preimage for each element of the target quotient. If a target
sequence is bad, its sequence of chosen preimages is bad, because any
comparison in the source would project to a comparison in the target.
This embeds the target bad-sequence tree into the source tree. The same
argument shows the target is a well-quasi-order whenever the source is.

\emph{Second argument, on fibres.}
Choose a linear extension of the target $X$ of type $\mot(X)$ and order the
fibres of the surjection consecutively according to that extension, then
linearize each fibre arbitrarily. Because the map is order preserving, no
comparison of the source is violated by this arrangement: a source comparison
projects to a target comparison, which is respected by the chosen ordering of
fibres, and comparisons inside a single fibre are respected by its own
linearization. The result is a linear extension of $Y$ of type at least
$\mot(X)$, so $\mot(X)\le\mot(Y)$. The first argument also yields the
well-quasi-order assertion; the second exhibits a concrete linear extension
and is what the exact-length upper bounds of Section~\ref{sec:exact} reuse.

\subsection{The ordinal arithmetic used here}
For ordinals in Cantor normal form, natural sum $\oplus$ adds coefficients
at equal exponents, and natural product $\otimes$ multiplies the resulting
formal polynomials, using natural sum on exponents. In particular,
\[
 (\omega^\alpha)\otimes(\omega^\beta)=\omega^{\alpha\oplus\beta}.
\]
Consequently, for finite $N\ge1$ and $q\ge1$,
\begin{equation}\label{eq:power}
 H_N^{\otimes q}=\omega^{\omega^{N-1}q},
\end{equation}
where the multiplication by $q$ on the right is ordinary ordinal
multiplication by a positive integer. Continuity of ordinal exponentiation
gives
\begin{equation}\label{eq:sup}
 \sup_{q\ge1}H_N^{\otimes q}
 =\omega^{\sup_{q\ge1}\omega^{N-1}q}
 =\omega^{\omega^N}=H_{N+1}.
\end{equation}
Equation~\eqref{eq:sup} is the entire ordinal calculation behind the
main induction.

A companion identity, using \emph{ordinary} ordinal exponentiation only, will
be needed for the alternative route of Section~\ref{sec:altroute}: for
$q\ge1$,
\begin{equation}\label{eq:amplification}
 \sup_{r\ge1}\Hn q^{\,r}
  =\sup_{r\ge1}\omega^{\omega^{\,q-1}\cdot r}
  =\omega^{\omega^{q}}
  =\Hn{q+1}.
\end{equation}

Ordinary powers and natural powers must not be interchanged in general, and
\eqref{eq:power}--\eqref{eq:sup} and \eqref{eq:amplification} are emphatically
not typographical variants of one another. In \eqref{eq:power} the exponent
$\otimes q$ is a natural power of $\Hn N$, computed by natural multiplication
of Cantor normal forms; in \eqref{eq:amplification} the exponent $r$ is an
ordinary ordinal power, computed by iterated ordinary multiplication. The two
displays happen to have the same supremum, which is exactly why the two routes
to the lower bound reach the same ordinal, and each display states exactly the
special case that its own route uses. Neither may be substituted for the other
inside a proof.

\section{Periodic tails and the finite generating alphabet}
\label{sec:normalform}

\subsection{One periodic block for each nonempty downset}
For a nonempty downset $D=\{d_0,\ldots,d_{t-1}\}$, fix a periodic word
\begin{equation}\label{eq:block}
 B_D=(d_0\cdots d_{t-1})^\omega.
\end{equation}
The chosen enumeration is immaterial up to $\eqv$: every element of $D$
appears infinitely often in every such representative.

\begin{lemma}[Comparison of periodic blocks]\label{lem:blocks}
For nonempty downsets $D,E$ and a letter $a\in P$,
\[
 a\emb B_D\ \Longleftrightarrow\ a\in D,
 \qquad
 B_D\emb B_E\ \Longleftrightarrow\ D\subseteq E.
\]
Every nonempty tail of $B_D$ is equivalent to $B_D$, and a letter $a$ embeds
into every tail of $B_D$ exactly when $a\in D$. Consequently, if every letter
of a finite word $p$ lies in $D$, then
\begin{equation}\label{eq:blockabsorb}
 pB_D\eqv B_D.
\end{equation}
\end{lemma}
\begin{proof}
If $a\le_P d$ for a letter $d\in D$, downward closure gives $a\in D$.
Conversely, $a\in D$ supplies a matching occurrence in $B_D$. If
$B_D\emb B_E$, applying the same observation to an occurrence of each
$d\in D$ gives $D\subseteq E$. If $D\subseteq E$, recursively match the
letters of $B_D$ to successive identical occurrences in $B_E$. There is
always a later occurrence, so the recursion gives an embedding.
The identical recursive construction works in any nonempty tail of a
periodic block; the reverse embedding into the full block is inclusion.
For \eqref{eq:blockabsorb}, place the finitely many letters of $p$ in an
initial portion of $B_D$, which is possible because each of them lies in $D$
and every element of $D$ recurs, and embed the remaining copy of $B_D$ into
the tail that is left. The reverse embedding simply omits $p$.
\end{proof}

Note that \eqref{eq:blockabsorb} holds for an arbitrary nonempty downset, not
only for $D=P$. The special case $D=P$, where $p$ may be any finite word at
all, is isolated later as \eqref{eq:absorb}.

For a nonempty subset $R\subseteq P$, repeating $R$ or repeating its downset
$\down R$ gives equivalent $\omega$-words. In the less immediate direction,
each $x\in\down R$ can be matched to an element $r\in R$ above it, and each
such $r$ recurs infinitely often.

\begin{remark}[Repeating only the maximal elements]\label{rem:maxblock}
Let $\Max D$ denote the set of maximal elements of a nonempty downset $D$.
Applying the previous paragraph with $R=\Max D$, whose downset is $D$, gives
\[
 B_{\Max D}\eqv B_D .
\]
Some presentations of this material define the periodic representative by
repeating $\Max D$ rather than all of $D$. The two conventions are
interchangeable everywhere below, and we keep the all-elements block $B_D$ of
\eqref{eq:block} as the official definition. Where a proof below inspects the
letters of a block, it only ever uses that the letters lie in $D$ and that
every element of $D$ recurs, which both representatives satisfy.
\end{remark}

\subsection{Reduction of arbitrary words}
For a word $u:\omega\to P$ put
\begin{equation}\label{eq:recdef}
 \rec(u)=\{a\in P:a\text{ occurs infinitely often in }u\},
 \qquad
 I(u)=\down\rec(u).
\end{equation}
The set $\rec(u)$ is nonempty because $P$ is finite and $u$ is infinite. Its
downward closure $I(u)$, and not the literal recurrent set, is the correct
invariant here: an embedding may increase labels, so what a block can absorb
is determined by the downset generated by its recurrent letters. We call
$I(u)$ the \emph{recurrent downset} of $u$.

The following periodic-tail reduction is standard; compare
\cite[Propositions~3.3--3.4]{Altman2026}. A proof is included to specify
exactly which equivalence is needed.

\begin{lemma}[Finite prefix and recurrent tail]\label{lem:recurrent}
Every word $u:\omega\to P$ is equivalent to $vB_D$ for some finite word
$v\in P^*$ and some nonempty downset $D$; one may take $D=I(u)$.
\end{lemma}
\begin{proof}
Let $R=\rec(u)$ be the set of labels occurring infinitely often in $u$. Since $P$
is finite and $u$ is infinite, $R$ is nonempty. Every label outside $R$
occurs only finitely often. Choose an index after all such occurrences,
and write $u=vt$, where $v$ is finite and $t$ uses only labels from $R$.
Every element of $R$ occurs infinitely often in $t$.

Put $D=\down R$. To embed $t$ into $B_D$, successively choose later identical
occurrences of its labels. To embed $B_D$ into $t$, for each requested label
$d\in D$ choose an $r\in R$ with $d\le_P r$, and then choose an occurrence
of $r$ later than all previous choices. Infinitely many occurrences of $r$
make this possible. Thus $t\eqv B_D$, and monotonicity of concatenation
gives $u\eqv vB_D$.
\end{proof}

The prefix can be normalized further, and the strengthened form is what the
canonical representatives of Section~\ref{sec:canonical} and the exact-length
computations of Section~\ref{sec:exact} require.

\begin{lemma}[One-block normal form]\label{lem:oneblock}
Every word $u:\omega\to P$ is equivalent to $pB_D$ with $D=I(u)$, where the
finite word $p$ may in addition be required to be empty or to end in a letter
outside $D$.
\end{lemma}
\begin{proof}
Each letter outside $D$ occurs only finitely often in $u$, and there are
finitely many letters, so either no letter outside $D$ occurs, in which case
put $p=\eps$, or there is a last such occurrence, in which case let $p$ be the
prefix ending at that occurrence. In both cases write $u=pt$ with every label
of $t$ inside $D$.

Every element of $D$ still recurs in $t$: given $d\in D$ there is, by
$D=\down\rec(u)$, a recurrent $b\ge_P d$, and only finitely many positions
were removed with $p$, so $b$ still recurs in $t$. Greedy matching of the
labels of $t$ to dominating occurrences gives $t\emb B_D$, and matching each
requested $d\in D$ to a later occurrence of such a $b$ gives $B_D\emb t$.
Hence $t\eqv B_D$, and concatenation gives $u\eqv pB_D$.
\end{proof}

\begin{proposition}[Finite representation]\label{prop:representation}
Every word of length less than $\omega^2$ is equivalent to a finite
concatenation of one-letter words and blocks $B_D$, with
$\varnothing\ne D\in\J(P)$. More precisely, a word of length $\omega r+k$ is
equivalent to an expression
\begin{equation}\label{eq:normal}
 p_0B_{D_0}p_1B_{D_1}\cdots p_{r-1}B_{D_{r-1}}p_r,
\end{equation}
in which every $p_i$ is a finite word, every $p_i$ with $i<r$ is empty or ends
in a letter outside $D_i$, and $p_r$ has length $k$.
\end{proposition}
\begin{proof}
Split a word of length $\omega r+k$ into its $r$ consecutive canonical
intervals of length $\omega$ and its final finite interval of length $k$.
Apply Lemma~\ref{lem:recurrent} to each infinite interval and concatenate the
resulting representatives; applying Lemma~\ref{lem:oneblock} instead yields
the normalized prefixes of \eqref{eq:normal}.
\end{proof}

By itself this is a representation theorem, not a uniqueness theorem: without
the extra clause in \eqref{eq:normal}, different finite expressions represent
equivalent words, for example $aB_{\down a}\eqv B_{\down a}$. Such absorption
is precisely why a finite generating alphabet gives only an upper bound
without further work. Uniqueness is nevertheless recoverable once the prefixes
are normalized as in \eqref{eq:normal}; this is carried out, with an explicit
stack-based procedure, in Section~\ref{sec:canonical}.

The construction is existential for an arbitrary infinite input word, since it
uses knowledge of which letters recur. The algorithm of
Section~\ref{sec:algorithm} takes finite expressions as input; it is not an
algorithm for discovering recurrence in an arbitrary stream.

\begin{definition}[Selected families]\label{def:selected}
For $\F\subseteq\J(P)\setminus\{\varnothing\}$, let $\W(P,\F)$ be the
quasi-order of all finite concatenations of letters in $P$ and periodic
blocks $B_D$ with $D\in\F$, using actual ordinal-word embeddability.
Finite expressions are identified up to mutual embeddability when maximal
order type is taken.
\end{definition}

In particular, $\W(P,\varnothing)=P^*$, and
$\W(P,\J(P)\setminus\{\varnothing\})$ has the same quotient order as
$\Wordspace P$ by Proposition~\ref{prop:representation}.

\begin{remark}[An intrinsic description of $\W(P,\F)$]\label{rem:intrinsic}
Definition~\ref{def:selected} is syntactic, but the resulting class of words
has an intrinsic description: up to equivalence, $\W(P,\F)$ consists exactly
of those words below $\omega^2$ each of whose canonical $\omega$-blocks has
its recurrent downset in $\F$. A finite prefix cannot change the recurrent
downset of a block, which gives one inclusion, and
Lemma~\ref{lem:oneblock} gives the converse by rewriting each block as
$pB_D$ with $D=I(u)\in\F$.
\end{remark}

\begin{proposition}[Atom-count upper bound]\label{prop:upper}
For $n=|P|\ge1$ and $m=|\F|$,
\begin{equation}\label{eq:upper}
 \mot(\W(P,\F))\le H_{n+m}.
\end{equation}
Moreover, $\W(P,\F)$ is a well-quasi-order.
\end{proposition}
\begin{proof}
Take a discrete alphabet of $n+m$ symbols: one for each finite letter and
one for each selected periodic block. Evaluate a finite symbol string by
ordinal concatenation. If one string is a subsequence of another, the
corresponding evaluated word embeds into the other by
Lemma~\ref{lem:concat}. Evaluation is surjective onto the represented
quasi-order. Apply the finite-word formula and monotone-surjection
principle in Theorem~\ref{thm:classical}.
\end{proof}

Using a discrete alphabet deliberately ignores extra comparisons among the
atoms. This does not weaken the desired upper bound: every poset with the
same positive finite cardinality has the same finite-word maximal order
type. No order-reflection claim is made for the evaluation map.

\section{Cofinal localization and synchronization}
\label{sec:sync}

The following lemma concerns \emph{displayed target tokens}. A finite word
may contain many one-letter tokens adjacent to each other, and a token
expression need not be a canonical decomposition of its ordinal length.
Nevertheless, each displayed $B_E$ occupies a convex interval of order type
$\omega$, and there are only finitely many displayed tokens.

\begin{lemma}[Cofinal localization]\label{lem:cofinal}
Suppose a periodic block $B_D$ embeds into a finite token expression $v$.
Then a tail of its image lies in one displayed target block $B_E$, is
cofinal in that block, and satisfies $D\subseteq E$.
If a source position occurs after the whole source block $B_D$, its image
occurs after the whole target block $B_E$.
\end{lemma}
\begin{proof}
Number the target tokens in their order of occurrence. As the source
positions $0,1,2,\ldots$ increase, the indices of the tokens containing
their images form a nondecreasing sequence with only finitely many
possible values. It is eventually constant. The final token cannot be a
one-letter token, because the embedding is injective. It is therefore
some $B_E$. The image contains infinitely many positions of that token,
so it is unbounded in its order type $\omega$.

Every $d\in D$ still occurs in the source tail. Each such occurrence maps
to a label in $E$ above $d$. Downward closure gives $d\in E$, proving
$D\subseteq E$. Finally, the image of any subsequent source position must
lie above every image point of $B_D$. Since those points are cofinal in the
target $B_E$, no position of that target block remains available.
\end{proof}

\begin{remark}[A block may straddle earlier blocks]\label{rem:straddle}
The lemma does not say that the entire source block lands in a single
target block. An initial finite part can be placed earlier. Only the tail
is localized. All subsequent proofs use this weaker, correct statement.
Under the stronger hypothesis that source and target have the \emph{same}
number of canonical $\omega$-blocks, whole-block containment does hold; that
is Lemma~\ref{lem:equalblocks} below, which is used only in
Sections~\ref{sec:exact} and~\ref{sec:canonical}. The two statements are
compatible because their hypotheses differ, and nothing in
Lemma~\ref{lem:equalblocks} retracts the present remark.
\end{remark}

\begin{remark}[The role of explicit segments]\label{rem:explicit}
A displayed token $B_E$ inside an expression need not begin at a canonical
$\omega$-block boundary of the word it denotes: if a finite word precedes it,
its pure periodic segment starts at a finite offset inside the corresponding
canonical block. That segment still has length $\omega$, is cofinal in the
block, and ends at the next canonical boundary; this is the ordinal identity
$k+\omega=\omega$ again. Throughout the arguments below, ``inside a displayed
separator'' therefore always means inside that \emph{explicit pure periodic
segment}, never inside the finite prefix that happens to precede it. This is a
statement about the \emph{target} side, complementary to
Remark~\ref{rem:straddle}, which is about a \emph{source} block straddling
earlier target blocks.
\end{remark}

\begin{lemma}[Equal-count synchronization]\label{lem:sync}
Let $D$ be a nonempty downset, and suppose an old family $\G$ satisfies
\begin{equation}\label{eq:newness}
  D\nsubseteq E\qquad\text{for every }E\in\G.
\end{equation}
Consider two token expressions with exactly $r$ displayed copies of $B_D$,
all other infinite tokens belonging to $\G$. Under an embedding of the
first expression into the second, the tail of the $i$th displayed source
$B_D$ is cofinal in the $i$th displayed target $B_D$, for $1\le i\le r$.
\end{lemma}
\begin{proof}
By Lemma~\ref{lem:cofinal}, each displayed source $B_D$ requires a target
infinite block whose downset contains $D$. Condition~\eqref{eq:newness}
excludes all old target blocks. Its cofinal destination is therefore one
of the $r$ displayed new target blocks. Successive source new blocks have
strictly increasing cofinal destination indices: after the first has used
a target block cofinally, any later source position must be beyond it.
A strictly increasing injection from $\{1,\ldots,r\}$ into itself is the
identity. This proves the claim.
\end{proof}

\begin{corollary}[Location between synchronized boundaries]\label{cor:between}
In the situation of Lemma~\ref{lem:sync}, write the source expression as
$b_0B_Db_1B_D\cdots B_Db_r$ and the target as $d_0B_Dd_1B_D\cdots B_Dd_r$,
with all $b_i,d_i$ in $\W(P,\G)$. Then, for each $i$:
\begin{enumerate}[label=\textnormal{(\alph*)}]
\item if $i>0$, every image point of $b_i$ occurs at or after the end of the
      $i$th displayed target separator;
\item if $i<r$, every image point of $b_i$ occurs before the image of the
      first point of the $(i+1)$st source separator, hence before the end of
      the $(i+1)$st displayed target separator;
\item $b_r$ embeds into $d_r$.
\end{enumerate}
\end{corollary}
\begin{proof}
Every point of $b_i$ follows every point of the preceding source separator.
That separator's image is cofinal in the corresponding target separator by
Lemma~\ref{lem:sync}, so no target position of that separator remains
available, giving~(a). Every point of $b_i$ precedes the next source
separator, whose first image point lies before the end of the corresponding
target separator, again by cofinality; strict monotonicity gives~(b).
After the last target separator, the only remaining target material is $d_r$,
which with~(a) gives~(c).
\end{proof}

Part~(b) still permits a finite tail of $b_i$ to enter the next target
separator. The guards of Sections~\ref{sec:proper} and~\ref{sec:universal}
eliminate precisely that possibility.

To arrange \eqref{eq:newness}, introduce a finite collection of downsets in
an inclusion-respecting order. Ordering them by cardinality, with arbitrary
tie breaking, suffices: a new set cannot be contained in an earlier,
distinct set of no larger size. The full downset $P$, when present, is last.

\section{Proper markers: finite guards give product embeddings}
\label{sec:proper}

\begin{lemma}[Right cancellation of a final letter]\label{lem:cancel}
For arbitrary ordinal words $u,v$ and any fixed letter $c$,
\[
 uc\emb vc\quad\Longrightarrow\quad u\emb v.
\]
\end{lemma}
\begin{proof}
The image of the final source letter lies at or before the final target
letter. Every position belonging to $u$ maps strictly before that image,
hence inside the target prefix $v$. Restrict the embedding.
\end{proof}

\begin{theorem}[Proper-marker product embedding]\label{thm:proper}
Let $D$ be a nonempty proper downset of $P$, and let $\G$ satisfy
\eqref{eq:newness}. Choose $c\in P\setminus D$ and put
$X=\W(P,\G)$, $Y=\W(P,\G\cup\{D\})$. For every $r\ge1$, the map
\begin{equation}\label{eq:propermap}
 \Phi_r(u_0,\ldots,u_r)
   =u_0cB_Du_1cB_D\cdots u_{r-1}cB_Du_r
\end{equation}
is an order embedding $X^{r+1}\hookrightarrow Y$.
\end{theorem}
\begin{proof}
Order preservation follows by concatenating componentwise embeddings and
identity maps on the separators. It remains to prove reflection.

Suppose $\Phi_r(u_0,\ldots,u_r)\emb\Phi_r(v_0,\ldots,v_r)$.
By Lemma~\ref{lem:sync}, the $i$th source separator $B_D$ has image tail
cofinal in the $i$th target separator $B_D$. Thus, for $i>0$, every image
point of $u_i$ and of its subsequent guard lies after the $i$th target
separator. For $i=0$ there is no preceding separator and no lower-bound
restriction to impose.

Fix $i<r$. The source guard $c$ immediately following $u_i$ precedes the
next source separator. Some point of that next separator maps into the
corresponding next target separator. Therefore the source guard's image
is before the end of that target separator. It cannot be in that separator:
if it were matched to a label $d\in D$, then $c\le_P d$ would imply
$c\in D$, contrary to the choice of $c$. Its image consequently lies in
the target interval $v_ic$ before the separator. Every point of $u_i$
precedes the guard, so its image lies in $v_i$, not in the final guard of
$v_ic$. We obtain $u_i\emb v_i$.

After the last source separator, cofinality forces the entire source
suffix $u_r$ to map into the target suffix $v_r$. Hence $u_r\emb v_r$ as
well. All components compare in the product order. Reflection and
preservation imply an order embedding after taking equivalence classes.
\end{proof}

The guard is allowed to map into the interior of $v_i$; it need not match
the explicitly displayed target guard. The proof only needs the fact that
it lies no later than that guard. This is important when old components
already contain many occurrences of $c$ or old infinite blocks.

\begin{corollary}[One proper marker raises the level]\label{cor:properlevel}
Under the hypotheses of Theorem~\ref{thm:proper}, if
$\mot(X)=H_N$ and $n+|\G|=N$, then $\mot(Y)=H_{N+1}$.
\end{corollary}
\begin{proof}
For every $r\ge1$, product embedding and the classical product theorem give
$\mot(Y)\ge H_N^{\otimes(r+1)}$. Taking the supremum and using
\eqref{eq:sup} yields $\mot(Y)\ge H_{N+1}$. Proposition~\ref{prop:upper}
gives the reverse inequality.
\end{proof}

\section{The universal marker: limit padding replaces a guard}
\label{sec:universal}

For $D=P$ there is no excluded letter. Any fixed finite word $z$ is absorbed
by the universal periodic block:
\begin{equation}\label{eq:absorb}
 zB_P\eqv B_P.
\end{equation}
Indeed, the countable sequence $zB_P$ can be matched greedily into $B_P$;
every requested label reappears infinitely often. Thus attaching a finite
guard before $B_P$ cannot preserve arbitrary component information.
The remedy is to end the component with a cancellable \emph{infinite} tail.

\subsection{A cancellable map into limit lengths}

\begin{lemma}[Limit padding]\label{lem:padding}
Let $E$ be a nonempty proper downset of $P$, and choose $c\notin E$.
The map
\begin{equation}\label{eq:padding}
 G(u)=ucB_E
\end{equation}
is order preserving and order reflecting on arbitrary words of length
below $\omega^2$. Every $G(u)$ has nonzero limit length, of the form
$\omega k$ with $k\ge1$. If $B_E$ is allowed in a family, $G$ maps that
family's represented word order into itself.
\end{lemma}
\begin{proof}
Monotonicity is immediate. Suppose $ucB_E\emb vcB_E$. The source occurrence
of $c$ cannot map into the final target $B_E$, by downward closure and
$c\notin E$. Thus it maps into $vc$, and the restriction to $u$ maps into
$v$, by the same final-letter argument as Lemma~\ref{lem:cancel}. Hence
$u\emb v$.

If $\len(u)=\omega k+m$, then
\[
 \len(ucB_E)=(\omega k+m)+1+\omega=\omega(k+1),
\]
which is a nonzero limit ordinal. The family assertion follows from the
explicit concatenation.
\end{proof}

\begin{lemma}[No infinite tail fits below a finite bound]\label{lem:noleak}
Let $u$ have nonzero limit length below $\omega^2$. Suppose an embedding
maps all of $u$ to positions preceding a fixed position $\eta$ inside a
target interval of order type $\omega$. Then no position of $u$ can map
inside that interval.
\end{lemma}
\begin{proof}
If a source position $\xi$ mapped into the interval before $\eta$, there
would be only finitely many target positions between its image and
$\eta$. But a nonzero limit ordinal has infinitely many positions after
each of its positions: otherwise the ordinal would have a last position.
The infinitely many source positions after $\xi$ would have to inject into
that finite target interval. This is impossible.
\end{proof}

The position $\eta$ need not be the first position of the target
$\omega$-block. Any point of it gives a finite bound. Nor is it necessary
for the next source marker to begin inside that block: one image point of
its cofinal tail suffices.

The lemma must not be misread as saying that a word of limit length never
embeds into a universal periodic word. It certainly may; $B_P$ alone embeds
into $B_P$. The obstruction exploited below is different, and it is created by
the surrounding configuration: a further synchronized source separator is
required to begin at an actual target position before the end of that same
target block, so the padded component and that separator would have to share
one interval of order type $\omega$ with only finitely many positions between
them.

\subsection{The product embedding}

\begin{theorem}[Universal-marker product embedding]\label{thm:full}
Let $\G$ consist of nonempty proper downsets of $P$, with
$\G\ne\varnothing$. Choose $E\in\G$ and the map $G$ from
Lemma~\ref{lem:padding}. Put
$X=\W(P,\G)$ and $Y=\W(P,\G\cup\{P\})$. For each $r\ge1$,
\begin{equation}\label{eq:fullmap}
 \Psi_r(u_0,\ldots,u_r)
  =G(u_0)B_PG(u_1)B_P\cdots G(u_{r-1})B_Pu_r
\end{equation}
is an order embedding $X^{r+1}\hookrightarrow Y$.
\end{theorem}
\begin{proof}
Preservation follows from the monotonicity of $G$ and concatenation.
Suppose the image of $(u_0,\ldots,u_r)$ embeds into the image of
$(v_0,\ldots,v_r)$.

No old downset contains $P$. Lemma~\ref{lem:sync} therefore applies to the
$r$ displayed universal separators. Their cofinal destinations are the
corresponding $r$ displayed target universal separators.

Fix $i<r$. When $i>0$, cofinality at the preceding separator forces all of
$G(u_i)$ to map after the corresponding preceding target separator.
Choose an image point $\eta$ of the following source $B_P$ inside the
corresponding following target $B_P$; such a point exists by cofinal
localization. Every image point of $G(u_i)$ is before $\eta$.
By Lemma~\ref{lem:padding}, $G(u_i)$ has nonzero limit length. Applying
Lemma~\ref{lem:noleak} shows that none of its positions can be in this
following target $B_P$. Thus the whole image of $G(u_i)$ lies in the
intervening target component $G(v_i)$. For $i=0$ the same argument applies
without a preceding boundary. We have proved
\[
 G(u_i)\emb G(v_i)\qquad(i<r).
\]
Reflection for $G$ yields $u_i\emb v_i$ for all $i<r$.

After the last source universal separator, cofinality forces the final
source component $u_r$ into the final target component $v_r$. Hence
$u_r\emb v_r$, completing product reflection.
\end{proof}

\begin{corollary}[One universal marker raises the level when padding exists]
\label{cor:fulllevel}
Under the hypotheses of Theorem~\ref{thm:full}, if
$\mot(X)=H_N$ and $n+|\G|=N$, then $\mot(Y)=H_{N+1}$.
\end{corollary}
\begin{proof}
Apply the finite Cartesian-product theorem to the embeddings $\Psi_r$,
take the supremum in \eqref{eq:sup}, and combine with the atom-count
upper bound, exactly as in Corollary~\ref{cor:properlevel}.
\end{proof}

Notice the role of the old proper block. It need not encode the whole
component by itself. Its job is to give a limit-length ending, while an
excluded finite letter makes that ending cancellable. These two properties
are logically distinct and both are used.

\subsection{The two constructions as one statement}
\label{sec:unified}

Theorems~\ref{thm:proper} and~\ref{thm:full} differ only in which guard map is
attached to a component. It is convenient, and it is what the quick reference
of Appendix~\ref{app:quickref} states, to record them as a single theorem.

\begin{theorem}[Unified guarded product embedding]\label{thm:product}
Let $D$ be a nonempty downset not in $\G$ with $D\nsubseteq E$ for every
$E\in\G$, and assume either
\begin{enumerate}[label=\textnormal{(\roman*)}]
\item $D\ne P$, or
\item $D=P$ and $\G$ contains a nonempty proper downset.
\end{enumerate}
Define a guard map $g_D$ on $\W(P,\G)$ by
\begin{equation}\label{eq:unifiedguard}
 g_D(u)=
 \begin{cases}
  uc,&D\ne P,\ c\in P\setminus D,\\
  ucB_R,&D=P,\ R\in\G\text{ proper},\ c\in P\setminus R,
 \end{cases}
\end{equation}
and put, for $m\ge0$,
\begin{equation}\label{eq:unifiedmap}
 \Phi_m(u_0,\ldots,u_m)=g_D(u_0)B_Dg_D(u_1)B_D\cdots B_Dg_D(u_m).
\end{equation}
Then
\[
 \Phi_m(u_0,\ldots,u_m)\emb\Phi_m(v_0,\ldots,v_m)
 \iff u_i\emb v_i\ \text{for every }i\le m,
\]
so $\Phi_m$ is an order embedding
$\W(P,\G)^{m+1}\hookrightarrow\W(P,\G\cup\{D\})$.
\end{theorem}
\begin{proof}
Both guard maps are order preserving and order reflecting and send
$\W(P,\G)$ into itself: for case~(i) this is Lemma~\ref{lem:cancel} together
with the concatenation lemma, and for case~(ii) it is
Lemma~\ref{lem:padding}. Preservation of $\Phi_m$ follows by concatenating
componentwise embeddings with identity maps on the separators.

For reflection, Lemma~\ref{lem:sync} aligns the $m$ displayed separators and
Corollary~\ref{cor:between} places each guarded component between the
corresponding target boundaries. In case~(i) the terminal letter $c$ of
$g_D(u_i)$ cannot lie inside the next explicit target separator, because
$c\notin D$ and $D$ is downward closed, and it cannot lie after that separator
by Corollary~\ref{cor:between}(b); so the whole of $g_D(u_i)$ lies in
$g_D(v_i)$. In case~(ii) the word $g_D(u_i)$ has nonzero limit length by
Lemma~\ref{lem:padding}, so Lemma~\ref{lem:noleak} forbids it from meeting the
next explicit target separator at all, with the same conclusion. The last
component is handled by Corollary~\ref{cor:between}(c). Reflection of $g_D$
then gives $u_i\emb v_i$ for every $i$.
\end{proof}

\begin{remark}[Guarding the final component or not]\label{rem:lastguard}
Theorem~\ref{thm:product} guards \emph{every} component, including the last
one, whereas the maps $\Phi_r$ of \eqref{eq:propermap} and $\Psi_r$ of
\eqref{eq:fullmap} leave the final component $u_r$ unguarded. Both conventions
are correct and both are used in the literature on this construction. With a
trailing guard the final component is handled by cancellation, exactly like
the others; without one it is handled by Corollary~\ref{cor:between}(c),
since nothing follows it in the target. We take the unguarded form as the
official construction of Sections~\ref{sec:proper} and~\ref{sec:universal},
because it is the shorter expression, and record the all-guarded variant here
and in Appendix~\ref{app:quickref} so that the two displays are not mistaken
for a discrepancy.
\end{remark}

\subsection{An elementary alternative to the product theorem}
\label{sec:altroute}

Corollaries~\ref{cor:properlevel} and~\ref{cor:fulllevel} pass from the
Cartesian-power embeddings to the ordinal lower bound through
Theorem~\ref{thm:classical}(iii), the de Jongh--Parikh natural-product
formula. That import can be avoided. The following lemma supplies, by hand,
the only consequence of it that the main theorem needs, and the route is
recorded here in full because it changes the hypotheses of the main theorem
rather than merely restating it.

\begin{lemma}[Finite powers amplify a maximal order type]\label{lem:power}
For every well-quasi-order $X$ and every finite $r\ge1$,
\begin{equation}\label{eq:lexlower}
 \mot(X^r)\ \ge\ \mot(X)^r,
\end{equation}
with \emph{ordinary} ordinal exponentiation on the right.
\end{lemma}
\begin{proof}
Let $L$ be a maximal linear extension of the quotient of $X$, of type
$\alpha=\mot(X)$. Order $r$-tuples by comparing their \emph{last unequal
coordinate}, using $L$ in that coordinate. This is a linear order, and it is a
well-order of type $\alpha^r$: the last coordinate indexes $\alpha$
consecutive copies of the corresponding order on $(r-1)$-tuples, so the type
satisfies $\alpha_r=\alpha_{r-1}\cdot\alpha$ with $\alpha_1=\alpha$.

It extends the componentwise order. Indeed, if one tuple is componentwise
below another and the two differ, then at every coordinate the first is below
or equal in $L$, and at the last coordinate where they differ it is strictly
earlier. Hence the last-unequal-coordinate order is a linear extension of
$X^r$, and \eqref{eq:lexlower} follows from the definition of $\mot$.
\end{proof}

\begin{corollary}[One new block raises one Higman level, elementarily]
\label{cor:altlevel}
Let $X=\W(P,\G)$ and $Y=\W(P,\G\cup\{D\})$ be as in
Theorem~\ref{thm:product}, and suppose $\mot(X)=\Hn N$ with $n+|\G|=N$. Then
$\mot(Y)=\Hn{N+1}$.
\end{corollary}
\begin{proof}
For every $r\ge1$, Theorem~\ref{thm:product} gives an order embedding
$X^{r}\hookrightarrow Y$, so $\mot(Y)\ge\mot(X^{r})\ge\Hn N^{\,r}$ by
Theorem~\ref{thm:classical}(i) and Lemma~\ref{lem:power}. Taking the supremum
over $r$ and applying the ordinary-exponentiation identity
\eqref{eq:amplification} gives $\mot(Y)\ge\Hn{N+1}$.
Proposition~\ref{prop:upper}, with one more atom, gives the reverse
inequality.
\end{proof}

Only an explicit linear extension of a Cartesian product was used here. In
particular, the argument does not require any general formula for the maximal
order type of a product, nor any formula for lexicographic products of
arbitrary partial orders. This is what the alternative route buys: with
Corollary~\ref{cor:altlevel} in place of
Corollaries~\ref{cor:properlevel} and~\ref{cor:fulllevel}, the proof of
Theorem~\ref{thm:main} uses from Theorem~\ref{thm:classical} only the
monotonicity assertions~(i) and~(ii) and the finite-word formula~(v), and in
particular no natural-product theorem. The main line of the article
nevertheless keeps the classical route, both because it is the shorter
derivation and because assertions~(iii) and~(iv) are genuinely needed later, in
the exact-length calculations of Section~\ref{sec:exact}, which the
elementary bound \eqref{eq:lexlower} does not supply.

\section{The full classification and the main theorem}
\label{sec:classification}

\begin{theorem}[Selected-tail classification]\label{thm:family}
Let $P$ be a nonempty finite poset of size $n$, and let
$\F\subseteq\J(P)\setminus\{\varnothing\}$.
\begin{enumerate}[label=\textnormal{(\roman*)}]
\item If $\F\ne\{P\}$, then
\begin{equation}\label{eq:family}
 \mot(\W(P,\F))=H_{n+|\F|}
 =\omega^{\omega^{\,n+|\F|-1}}.
\end{equation}
This includes the empty family.
\item If $\F=\{P\}$, then
\begin{equation}\label{eq:exception}
 \mot(\W(P,\{P\}))=H_n\cdot\omega
    =\omega^{\omega^{n-1}+1}.
\end{equation}
\end{enumerate}
\end{theorem}

\begin{proof}[Proof of part~\textnormal{(i)}]
For the empty family, \eqref{eq:family} is the finite-word formula.
Otherwise $\F\ne\{P\}$ implies that the family contains at least one
proper downset. List its proper downsets as $D_1,\ldots,D_t$ in increasing
cardinality, breaking ties arbitrarily. If $P\in\F$, append it last.

Start with $X_0=P^*$, of type $H_n$. At step $i\le t$, no earlier downset
contains $D_i$. Indeed, an earlier downset has no greater cardinality, and
equality together with containment would make the downsets equal.
Theorem~\ref{thm:proper} applies. Inductively,
Corollary~\ref{cor:properlevel}, or equally
Corollary~\ref{cor:altlevel} if one prefers the elementary route, gives
\[
 \mot\bigl(\W(P,\{D_1,\ldots,D_i\})\bigr)=H_{n+i}.
\]
If $P$ is absent from $\F$, this finishes the proof. If $P$ is present,
there is an old proper downset because $t\ge1$. Apply
Theorem~\ref{thm:full} and Corollary~\ref{cor:fulllevel} to the final step,
obtaining $H_{n+t+1}=H_{n+|\F|}$.
\end{proof}

\begin{proof}[Proof of part~\textnormal{(ii)}]
Every expression using only the universal infinite block can be written
\[
 u_0B_Pu_1B_P\cdots u_{k-1}B_Pv,
 \qquad u_i,v\in P^*.
\]
By \eqref{eq:absorb}, it is equivalent to $B_P^kv$. Here $B_P^k$ means
$k$ successive copies, not an infinite repetition. These normal forms
satisfy
\begin{equation}\label{eq:univorder}
 B_P^kv\emb B_P^\ell w
 \quad\Longleftrightarrow\quad
 \bigl(k<\ell\bigr)\ \text{or}\
 \bigl(k=\ell\text{ and }v\emb w\bigr).
\end{equation}
If $k<\ell$, send the $k$ source infinite blocks into the first $k$ target
blocks and the finite suffix into one extra universal target block.
If $k>\ell$, $k$ successive source infinite blocks require $k$ distinct
cofinal target infinite destinations by Lemma~\ref{lem:cofinal}, which is
impossible. If $k=\ell$, synchronization leaves the source suffix wholly
in the target suffix, so $v\emb w$ is necessary. It is plainly sufficient.

Thus the quotient is the ordered sum of $\omega$ copies of $P^*$:
\begin{equation}\label{eq:sum}
 \W(P,\{P\})/\eqv\ \cong\ \sum_{k<\omega}P^*.
\end{equation}
Every linear extension must put all of the $k$th copy before the
$(k+1)$st copy. Each copy has maximal order type $H_n$, and one may use
a maximal linear extension in every copy. Therefore the maximal order
type is the ordinary ordered sum $\sum_{k<\omega}H_n=H_n\cdot\omega$.
The elementary identity
$\omega^\alpha\cdot\omega=\omega^{\alpha+1}$ gives
\eqref{eq:exception}.
\end{proof}

\begin{proof}[Proof of Theorem~\ref{thm:main}]
Suppose $n\ge2$. Choose a minimal element $a$ of $P$. The singleton
$\{a\}$ is a nonempty proper downset. Consequently the full family
$\J(P)\setminus\{\varnothing\}$ is not the exceptional family $\{P\}$.
Its size is $j(P)-1$. Proposition~\ref{prop:representation} and
Theorem~\ref{thm:family}(i) give
\[
 \mot\bigl(\Wordspace P\bigr)
 =H_{n+j(P)-1}=\omega^{\omega^{n+j(P)-2}}.
\]
For $n=1$, there is one word at each ordinal length $\alpha<\omega^2$, and
embedding is exactly comparison of lengths. Its quotient is therefore the
ordinal $\omega^2$ itself. Equivalently, apply part~(ii) with $H_1=\omega$.
For $n=0$, the empty word is the sole word, giving maximal order type $1$.
\end{proof}

\subsection{Why the supremum argument is legitimate}
\label{sec:supremum}
The proof never identifies the order type of an increasing union with the
supremum of the order types of its pieces. Such an identification is not a
general principle for maximal order types. Instead, for each fixed finite
$r$, there is a separately verified order embedding of a Cartesian power.
It gives a genuine lower bound on the \emph{same} target order type.
Taking the supremum of a collection of lower bounds is valid. The upper
bound is independently obtained by a monotone surjection.

The distinction is not pedantic, and Section~\ref{sec:exact} makes it
concrete. There the maximal order type of each individual length stratum is
computed exactly, and the supremum of those values, recorded in
\eqref{eq:stratumsup}, is $\omega^{\omega^n}$: strictly smaller than the
answer $\omega^{\omega^{n+j(P)-2}}$ of Theorem~\ref{thm:main} whenever
$n\ge2$. A stratum-by-stratum supremum therefore genuinely fails to compute
the theorem, while the supremum of lower bounds used here does not.

\subsection{The singleton warning is structural}
\label{sec:singleton}
For a singleton alphabet, two formal atoms are available: the letter and
its periodic $\omega$-block. Their finite words would have maximal type
$H_2=\omega^\omega$ if treated as a free Higman alphabet. But after ordinal
evaluation, the actual type is only $\omega^2$. The family
$\F=\{P\}$ exhibits the same defect over every nonempty finite alphabet.
Hence atom counting alone cannot be a lower-bound proof. The padding
construction explains exactly why the full family over $n\ge2$ avoids this
collapse.

\section{The binary proof written out explicitly}
\label{sec:binary}

This section specializes the general proof so the selected problem can be
checked without tracking an arbitrary downset family.

\subsection{Two incomparable letters}
Let $P=\{a,b\}$ with no comparisons between distinct letters, and abbreviate
\[
 A=a^\omega,\qquad B=b^\omega,\qquad U=(ab)^\omega.
\]
There are exactly three nonempty downsets, giving these three periodic
block types. Define stages by their allowed tokens, and record the guard used
to activate the block that is new in each row:
\[
 \begin{array}{c|c|c|c}
  \text{stage}&\text{allowed tokens}&\text{guard for the new block}
    &\text{claimed maximal type}\\\hline
  X_0&a,b&\text{none; finite words}&H_2=\omega^{\omega}\\
  X_1&a,b,A&u\mapsto ub&H_3=\omega^{\omega^2}\\
  X_2&a,b,A,B&u\mapsto ua&H_4=\omega^{\omega^3}\\
  X_3&a,b,A,B,U&u\mapsto ub\,a^\omega&H_5=\omega^{\omega^4}.
 \end{array}
\]
The upper bound in each row comes from the number of tokens. The guard in the
last row is the limit padding $G$ of Lemma~\ref{lem:padding}; the guards in
the two middle rows are single excluded letters.
The first equality is classical. To pass from $X_0$ to $X_1$, use, for
every finite $r$,
\[
 (u_0,\ldots,u_r)\longmapsto
 u_0bAu_1bA\cdots u_{r-1}bAu_r.
\]
Source $A$-blocks must synchronize with target $A$-blocks, and the guard
$b$ cannot enter $A$. This embeds $X_0^{r+1}$.

To pass from $X_1$ to $X_2$, use
\[
 (u_0,\ldots,u_r)\longmapsto
 u_0aBu_1aB\cdots u_{r-1}aBu_r.
\]
A source $B$-block cannot be cofinal in an old target $A$-block, and its
preceding guard $a$ cannot enter a target $B$-block. This embeds
$X_1^{r+1}$.

For the final step, set
\[
 G(u)=ubA.
\]
The guard $b$ cannot enter the final $A$, so $G$ reflects embedding; the
image has limit length. Now use
\[
 (u_0,\ldots,u_r)\longmapsto
 G(u_0)U G(u_1)U\cdots G(u_{r-1})Uu_r.
\]
A source $U$ contains both letters infinitely often, so neither $A$ nor
$B$ can be its cofinal destination. The displayed $U$-blocks synchronize.
Limit padding prevents the intervening components from entering the next
$U$ before a fixed image point of its source marker. This embeds
$X_2^{r+1}$.

At each of the three steps, \eqref{eq:sup} converts these product embeddings
into the next displayed ordinal. Finally, every binary word of length
below $\omega^2$ is equivalent to an expression in $a,b,A,B,U$.
This proves the sharpness of the binary antichain bound.

\subsection{Two ordered letters}
Let $0<1$. The nonempty downsets are $\{0\}$ and $\{0,1\}$; their periodic
blocks may be represented by
\[
 A=0^\omega,\qquad U=1^\omega.
\]
The latter is equivalent to $(01)^\omega$. Begin with finite words over
$\{0,1\}$, of maximal type $H_2$. Add $A$ using the guard $1$, giving
$H_3$. Then add $U$ using the limit padding $G(u)=u1A$, giving $H_4$.
Thus the chain's sharpened upper bound is also attained.

\subsection{Three easily confused comparisons}
For the antichain, the following facts illustrate the mechanism:
\[
 aU\eqv U,\qquad A\emb U,\qquad AB\not\emb U.
\]
The first is finite-prefix absorption. The second uses cofinal recurrence
of $a$ in $U$. The third follows because two successive $\omega$-blocks
cannot both have cofinal image in the single target block. Also
$U\not\emb AB$: its tail would have to lie in either the pure $a$ block
or the pure $b$ block, losing one of its recurrent labels. These facts
show why both order of occurrence and cofinal type matter.

\subsection{Negative controls: three invalid shortcuts}
\label{sec:negative}

The hypotheses of Theorems~\ref{thm:proper} and~\ref{thm:full} are not
decorative. Each of the following three simplifications is tempting, and each
fails on the two-letter antichain $P=\{a,b\}$ with an explicit witness. The
three witnesses are also executed as negative controls in the verification
suite of Section~\ref{sec:tests}, where they are required to \emph{fail},
so a regression that silently weakened a hypothesis would be caught.

\paragraph{Shortcut 1: omitting the guards.}
For any letter $a$,
\begin{equation}\label{eq:neg1}
 aB_P\eqv B_P,\qquad\text{while}\qquad a\not\emb\eps .
\end{equation}
A component placed directly before a separator can therefore disappear under
embeddability, so an unguarded tuple map need not reflect the componentwise
order. This is \eqref{eq:absorb} read as a counterexample.

\paragraph{Shortcut 2: using a finite guard at the universal block.}
Appending one letter does not repair the previous failure:
\begin{equation}\label{eq:neg2}
 aaB_P\emb aB_P,\qquad\text{while}\qquad a\not\emb\eps .
\end{equation}
More generally, any fixed finite padding is absorbed into a universal
periodic tail by \eqref{eq:absorb}. The guard of Lemma~\ref{lem:padding} is of
limit length for a mathematical reason, not for convenience.

\paragraph{Shortcut 3: inserting a downset after one that contains it.}
The newness condition \eqref{eq:newness} cannot be dropped. Let the old family
be $\G=\{E\}$ with $E=P$, and let the new downset be $D=\{a\}$, which is
contained in $E$; take the naive terminal guard $b\notin D$. Consider the two
component pairs
\[
 (\eps,\,a)\qquad\text{and}\qquad(B_E,\,\eps).
\]
The first is \emph{not} componentwise below the second, since
$a\not\emb\eps$. Nevertheless their guarded encodings satisfy
\begin{equation}\label{eq:negative}
 b\,a^\omega\,a\,b\ \emb\ B_E\,b\,a^\omega\,b .
\end{equation}
Indeed, match the initial source $b$ inside $B_E$, then the source $a^\omega$
cofinally inside that same $B_E$; then match the remaining source $a$ into the
displayed target $a^\omega$ and the final source $b$ to the terminal target
$b$. The source separator has used the \emph{old containing} block rather than
its designated target separator, and synchronization fails. Introducing
downsets in nondecreasing cardinality, as in the proof of
Theorem~\ref{thm:family}, is exactly what rules this out: a strictly earlier
downset of no larger size cannot contain the new one.

The article asserts the necessity of $D\nsubseteq E$ in prose several times;
\eqref{eq:negative} is the worked instance behind those assertions.

\section{Consequences and finite-poset examples}
\label{sec:consequences}

\subsection{Chains, antichains, and extremal values}

\begin{corollary}\label{cor:extrema}
For $n\ge2$,
\[
 \mot\bigl(\Wordspace {C_n}\bigr)=\omega^{\omega^{2n-1}},
 \qquad
 \mot\bigl(\Wordspace {A_n}\bigr)=\omega^{\omega^{n+2^n-2}}.
\]
Among $n$-element posets, the first value is the minimum, attained exactly
by chains, and the second is the maximum, attained exactly by antichains.
\end{corollary}
\begin{proof}
A chain has exactly $n+1$ downsets, its initial segments. Every subset of
an antichain is a downset, so it has $2^n$ downsets. For any finite poset,
the initial segments of a linear extension give $n+1$ distinct downsets,
and there are at most $2^n$ subsets in total.

For the equality statement at the lower end, fix a linear extension. If
there are only $n+1$ downsets, every downset is one of its initial segments.
In particular all principal downsets are comparable by inclusion. Since
$\down x\subseteq\down y$ implies $x\le_P y$, every pair of elements is
comparable, and $P$ is a chain. At the upper end, if every subset is a
downset, a strict relation $x<_P y$ would make $\{y\}$ fail to be a downset.
Thus $P$ is an antichain. The function of $j(P)$ in \eqref{eq:main} is
strictly increasing, so the ordinal extremal assertions follow.
\end{proof}

\subsection{Equivalent counting invariants}
Every downset of a finite poset is generated by its maximal elements,
which form an antichain. Conversely, the maximal elements of the downset
generated by an antichain are that antichain. These maps give a bijection,
including the empty set. Thus $j(P)$ may equally be described as the number
of antichains of $P$.

Complementation gives a bijection between downsets and upsets. Passing to
the dual poset exchanges downsets with upsets, so $j(P^{\mathrm{op}})=j(P)$.
The formula therefore yields
\[
 \mot\bigl(\Wordspace {P^{\mathrm{op}}}\bigr)
 =\mot\bigl(\Wordspace P\bigr).
\]
This is equality of maximal order types, not a claim that the word
quasi-orders are isomorphic.

For a disjoint union $P\amalg Q$, choosing a downset is equivalent to
choosing one downset in each component, hence
$j(P\amalg Q)=j(P)j(Q)$. For an ordered sum $P+Q$, in which every element
of $P$ lies below every element of $Q$, one has
$j(P+Q)=j(P)+j(Q)-1$. The two possibilities are a downset contained in $P$,
or all of $P$ together with a nonempty downset of $Q$. These identities
make \eqref{eq:main} explicit for many recursively assembled alphabets.

\subsection{A table of examples}
In Table~\ref{tab:examples}, $\rho(P)=n+j(P)-2$ is the finite exponent
parameter in $\omega^{\omega^{\rho(P)}}$.

\begin{table}[htbp]
\centering
\begin{tabular}{@{}lrrrl@{}}
\toprule
Alphabet $P$ & $n$ & $j(P)$ & $\rho(P)$ & Maximal order type\\
\midrule
Two-element chain &2&3&3&$\omega^{\omega^3}$\\
Two-element antichain &2&4&4&$\omega^{\omega^4}$\\
Three-element chain &3&4&5&$\omega^{\omega^5}$\\
Three-element $V$ or its dual &3&5&6&$\omega^{\omega^6}$\\
Two-element chain plus an isolated point &3&6&7&$\omega^{\omega^7}$\\
Three-element antichain &3&8&9&$\omega^{\omega^9}$\\
Four-element chain &4&5&7&$\omega^{\omega^7}$\\
Four-element diamond &4&6&8&$\omega^{\omega^8}$\\
Four-element antichain &4&16&18&$\omega^{\omega^{18}}$\\
Five-element antichain &5&32&35&$\omega^{\omega^{35}}$\\
\bottomrule
\end{tabular}
\caption{Exact consequences of Theorem~\ref{thm:main}. The diamond has one
bottom element, two incomparable middle elements, and one top element.}
\label{tab:examples}
\end{table}

The alphabet cardinality by itself is not enough: the two binary values
already differ. Conversely, nonisomorphic alphabets can have the same
maximal order type. For instance the three-element $V$ and its dual have
the same $n$ and $j$, while a four-element chain and a three-element chain
plus an isolated point have the same value of $n+j-2$ despite different
cardinalities.

\subsection{Finite quasi-orders}
For a finite quasi-order rather than a poset, first pass to its quotient
by mutual comparison. Projection of labels preserves and reflects word
embedding: a comparison of quotient labels is exactly the comparison of
any representatives. Therefore the theorem applies with $n$ equal to the
number of equivalence classes and $j$ equal to the number of downsets of
the quotient poset. Counting raw labels instead of classes can give an
incorrect answer.

\section{Exact-length strata}
\label{sec:exact}

This section computes the maximal order type of each individual length
stratum. The results are supplementary: they are not used anywhere in the
proof of Theorem~\ref{thm:main}, and a reader checking only the resolution of
the two-letter question may skip the section entirely. Their point is
negative, and it is made in Section~\ref{sec:stratumwarning}: the answer of
Theorem~\ref{thm:main} cannot be assembled stratum by stratum.

Unlike the main line, this section does use the natural-product and
natural-sum identities of Theorem~\ref{thm:classical}(iii) and~(iv). The
elementary substitute of Lemma~\ref{lem:power} is a one-sided bound and does
not suffice here.

Throughout, $P$ is a nonempty finite poset with $n=|P|$, and
\[
 \E_\alpha(P)=\{u:\len(u)=\alpha\},
 \qquad
 d=j(P)-2
\]
is the number of \emph{proper} nonempty downsets of $P$. Thus $d\ge1$ exactly
when $n\ge2$, and $d=0$ for a singleton alphabet.

\subsection{A single infinite block}

\begin{proposition}[Length exactly $\omega$]\label{prop:exactomega}
For every nonempty finite poset $P$,
\begin{align}
 \mot\bigl(\E_\omega(P)\bigr)&=\Hn n\cdot d+1,\label{eq:exactomega}\\
 \mot\bigl(s_{\omega+1}(P)\bigr)&=\Hn n\cdot(d+1)+1,\label{eq:atmostomega}
\end{align}
where $s_{\omega+1}(P)$ denotes the words of length at most $\omega$.
\end{proposition}
\begin{proof}
By Lemma~\ref{lem:oneblock}, every word of length $\omega$ is equivalent to
$pB_D$ with $D=I(u)$ and $p$ empty or ending outside $D$. Group these classes
into \emph{fibres} indexed by the nonempty downset $D$.

Fix a proper nonempty $D$ and consider its fibre. For two normalized forms,
\begin{equation}\label{eq:fibre}
 pB_D\emb qB_D\quad\Longleftrightarrow\quad p\emb q.
\end{equation}
The implication from right to left is monotonicity of concatenation. For the
other direction, the claim is trivial when $p=\eps$; otherwise the final
letter of $p$ lies outside $D$, so by Lemma~\ref{lem:blocks} it cannot be
matched inside the target periodic tail, and its image lies in the target
prefix $q$. Every earlier position of $p$ maps strictly before it, hence also
into $q$.

By \eqref{eq:fibre} the map $p\mapsto pB_D$ is an isomorphism from the
suborder of $P^*$ consisting of the admissible prefixes, that is the words
that are empty or end outside $D$, onto the fibre of $D$. Hence that fibre has
type at most $\mot(P^*)=\Hn n$. In the other direction choose $c\notin D$ and
map $u\in P^*$ to $ucB_D$, whose prefix $uc$ is admissible. This map is order
preserving, and order reflecting by \eqref{eq:fibre} together with
Lemma~\ref{lem:cancel}. Hence the fibre of $D$ has type exactly $\Hn n$.

The fibre of the full downset $P$ is a single equivalence class: by
\eqref{eq:absorb} every $pB_P$ is equivalent to $B_P$. That class lies above
every word of length $\omega$, since $D\subseteq P$ always holds.
Conversely, by Lemma~\ref{lem:cofinal}, an embedding between words of length
$\omega$ forces inclusion of their recurrent downsets, so comparisons across
fibres go only from smaller to larger downsets.

For the lower bound, order the $d$ proper nonempty downsets by nondecreasing
cardinality, concatenate maximal linear extensions of their fibres in that
order, and place the single class $B_P$ last. No comparison is violated,
because a comparison between distinct fibres always runs from a downset of
strictly smaller cardinality to a larger one. This is a linear extension of
type $\Hn n\cdot d+1$.

For the upper bound, erase every comparison between distinct fibres. The
result is a disjoint union of $d$ orders of type $\Hn n$ and one singleton,
whose maximal order type is
$\Hn n\oplus\cdots\oplus\Hn n\oplus1=\Hn n\cdot d+1$ by
Theorem~\ref{thm:classical}(iv); the natural sum of $d$ copies of
$\omega^{\omega^{n-1}}$ with $1$ collects the coefficients at one exponent and
leaves the finite part alone. Deleting comparisons cannot decrease the
maximal order type, so the original order has type at most this value. This
proves \eqref{eq:exactomega}.

For \eqref{eq:atmostomega}, adjoin the finite words. No infinite word embeds
into a finite word, so placing all of $P^*$ first respects every comparison
and adds one further layer of type $\Hn n$; erasing the comparisons between
$P^*$ and $\E_\omega(P)$ gives the matching upper bound, again by
Theorem~\ref{thm:classical}(iv).
\end{proof}

\subsection{Several blocks and a fixed finite tail}

\begin{lemma}[Rigidity at equal block count]\label{lem:equalblocks}
Let $u$ and $v$ each have exactly $r$ canonical $\omega$-blocks, and suppose
$u\emb v$. Then the $i$th canonical source block is mapped, entirely and
cofinally, inside the $i$th canonical target block, for every $i<r$, and the
finite source tail is mapped inside the finite target tail.
\end{lemma}
\begin{proof}
By Lemma~\ref{lem:cofinal}, applied to the periodic normal form of each
canonical source block, each source block has a cofinal destination among the
target blocks, distinct source blocks have distinct destinations, and the
assignment is increasing. An increasing injection of an $r$-element chain into
itself is the identity, so the $i$th source block is cofinal in the $i$th
target block.

The $i$th source block cannot have an image point before the end of the
$(i-1)$st target block, because the $(i-1)$st source block has already used
that block cofinally and every point of the $i$th source block comes later.
It cannot have an image point after the end of the $i$th target block either:
if it did, the infinitely many source positions following that one would map
beyond the $i$th target block, leaving only finitely many image points inside
it, contradicting cofinality there. Hence the whole $i$th source block lies
inside the $i$th target block. After the last block only the finite target
tail remains, which handles the tail.
\end{proof}

Lemma~\ref{lem:equalblocks} is strictly stronger than
Lemma~\ref{lem:cofinal} with Remark~\ref{rem:straddle}, and it is not
available in the main argument: it needs the hypothesis that source and target
have the same number of canonical blocks, which is exactly what the product
constructions of Sections~\ref{sec:proper} and~\ref{sec:universal} do not
assume. Within the present section that hypothesis is automatic, because the
length is fixed.

\begin{proposition}[Fixed-length product decomposition]\label{prop:exactproduct}
For all $r,k<\omega$ there is an isomorphism of quotient orders
\begin{equation}\label{eq:exactisom}
 \E_{\omega r+k}(P)\ \cong\ \E_\omega(P)^{\,r}\times P^{k},
\end{equation}
and consequently, writing $A=\Hn n\cdot d+1$,
\begin{equation}\label{eq:exacttype}
 \mot\bigl(\E_{\omega r+k}(P)\bigr)=A^{\otimes r}\otimes n^{k}.
\end{equation}
\end{proposition}
\begin{proof}
Concatenation is a bijection between $r$-tuples of length-$\omega$ words
together with a length-$k$ tail and words of length $\omega r+k$, since the
canonical block decomposition of such a word is unique. By
Lemma~\ref{lem:equalblocks}, an embedding between two such words restricts to
an embedding of the $i$th block into the $i$th block and of the tail into the
tail; conversely, componentwise embeddings concatenate by
Lemma~\ref{lem:concat}. A strictly increasing injection between finite chains
of the same length $k$ is the identity, so the induced relation on the tails
is exactly the componentwise label order of $P^k$. This proves
\eqref{eq:exactisom}, and \eqref{eq:exacttype} follows from
Theorem~\ref{thm:classical}(iii), Proposition~\ref{prop:exactomega}, and
$\mot(P^k)=n^k$ for the finite poset $P^k$.
\end{proof}

Because $\Hn n=\omega^{\omega^{n-1}}$, expanding the natural power in
\eqref{eq:exacttype} gives the Cantor normal form
\begin{equation}\label{eq:cnfexact}
 \mot\bigl(\E_{\omega r+k}(P)\bigr)
  =\sum_{i=r,r-1,\ldots,0}
     \omega^{\omega^{\,n-1}\cdot i}\Bigl(n^{k}\binom ri d^{\,i}\Bigr).
\end{equation}
The sum is written in decreasing $i$, so it is an ordinary Cantor-normal-form
sum, and terms with zero coefficient are omitted. The coefficient arises as
follows: in the expansion of $(\Hn n\cdot d+1)^{\otimes r}$ one chooses, from
each of the $r$ factors, either the leading monomial $\Hn n\cdot d$ or the
constant $1$; choosing the leading monomial in exactly $i$ of the factors can
be done in $\binom ri$ ways, contributes the finite coefficient $d^{\,i}$, and
produces the exponent $\omega^{n-1}\oplus\cdots\oplus\omega^{n-1}$ with $i$
summands, namely $\omega^{n-1}\cdot i$. The convention $d^{\,0}=1$ is in
force, including when $d=0$: for a singleton alphabet
\eqref{eq:cnfexact} correctly collapses to $n^k=1$.

\subsection{A warning about taking suprema over lengths}
\label{sec:stratumwarning}

Suppose $n\ge2$, so that $d\ge1$. Taking the supremum of \eqref{eq:cnfexact}
over all lengths gives
\begin{equation}\label{eq:stratumsup}
 \sup_{r,k<\omega}\mot\bigl(\E_{\omega r+k}(P)\bigr)=\omega^{\omega^{\,n}},
\end{equation}
because the leading exponents are $\omega^{n-1}\cdot r$, whose supremum over
$r$ is $\omega^{n}$, and finite coefficients and lower-order terms do not
affect that supremum.

The value \eqref{eq:stratumsup} is \emph{strictly below} the answer
$\omega^{\omega^{\,n+j(P)-2}}$ of Theorem~\ref{thm:main}. To see the
strictness, take a linear extension of $P$: its $n+1$ prefixes are distinct
downsets, so $j(P)\ge n+1$ and therefore $n+j(P)-2\ge 2n-1>n$ whenever
$n\ge2$.

There is no contradiction, and the explanation is the one already given in
Section~\ref{sec:supremum}. The length strata are \emph{not} ordinal-sum
layers of the whole order. Ordinal length is indeed monotone under
embeddability, by Lemma~\ref{lem:cofinal} for the block count and
Lemma~\ref{lem:equalblocks} for the tail. But an ordinal-sum layering would
require every word of a shorter length to embed into every word of a longer
one, and that fails at once: over the two-letter antichain, $b$ has length $1$
and $a^\omega$ has length $\omega$, yet $b\not\emb a^\omega$. A maximal linear
extension of the full order is therefore free to interleave strata and exploit
incomparabilities between them, and it does. The downset induction of
Section~\ref{sec:classification} detects that additional complexity; the
separate exact-length calculations, by construction, cannot.

\section{An exact symbolic embedding algorithm}
\label{sec:algorithm}

\subsection{Representation}
The supplied module \texttt{code/omega2.py} represents a word by a finite
sequence of tokens. A token \texttt{('f', x)} is a finite letter, and
\texttt{('w', D)} is $B_D$, with $D$ stored as a nonempty bitmask downset.
The poset relation is supplied by predecessor masks. The code validates
reflexivity, antisymmetry, and transitivity.

This representation does not claim to compute the recurrent set of an
arbitrary black-box infinite input. Rather, it decides embeddings between
already supplied finite periodic-block descriptions. The mathematical
normal-form theorem says such descriptions represent every equivalence
class in the target order.

\subsection{Greedy decision procedure}
Maintain a pointer to a target token, initially the first. Read source
tokens from left to right.

\begin{description}[style=nextline,leftmargin=0pt,labelwidth=0pt,labelindent=0pt]
\item[Source token is a finite letter $a$.]
Scan target tokens until finding either a finite letter $b$ with $a\le_P b$,
or an infinite token $B_E$ with $a\in E$. In the finite case consume the
letter and advance the pointer. In the infinite case leave the pointer at
$B_E$: after finitely many matches, an equivalent periodic tail remains.
\item[Source token is an infinite block $B_D$.]
Scan target tokens until finding an infinite token $B_E$ with $D\subseteq E$.
Consume its tail cofinally and advance the pointer to the following token.
\end{description}
If the required token does not exist, reject. Accept after every source
token has been processed.

The same rules, written as a table of the four possible source/target token
pairs, are the following. A letter token is written $a$ or $b$ and an
infinite token is written $[D]$ or $[E]$, denoting $B_D$ and $B_E$.

\begin{center}
\begin{tabular}{@{}llll@{}}
\toprule
Source & Target & Matching condition & Cursor after a match\\
\midrule
$a$ & $b$ & $a\le_P b$ & Advance past $b$\\
$a$ & $[E]$ & $a\in E$ & Remain at $[E]$\\
$[D]$ & $[E]$ & $D\subseteq E$ & Advance past $[E]$\\
$[D]$ & $b$ & Never & No match\\
\bottomrule
\end{tabular}
\end{center}

When the current target token cannot match, skip it and advance the cursor.
The last row is worth stating explicitly: an infinite source token never
matches a finite target token, since a strictly increasing map out of $\omega$
cannot land in a single position. The second row is the only one that leaves
the cursor where it is, because a finite match consumes only finitely many
positions of a periodic tail and leaves an equivalent tail behind.

\begin{theorem}[Exactness of the algorithm]\label{thm:algorithm}
The greedy procedure accepts exactly when the represented source ordinal
word embeds into the represented target ordinal word.
\end{theorem}
\begin{proof}
Within a periodic target block, removing any finite initial segment leaves
an equivalent tail. Thus the symbolic state need not record a phase or a
finite offset in that block.

For a finite source letter, the first eligible target token provides an
early possible match. Choosing a later token cannot leave a more useful
suffix. Within an infinite token, choosing an actual sufficiently early
eligible occurrence leaves an equivalent periodic tail. Therefore a
successful continuation after any possible match can also be made after
the greedy match.

For an infinite source block $B_D$, Lemma~\ref{lem:cofinal} shows that every
embedding must use some target block $B_E$ cofinally with $D\subseteq E$.
All later source positions must then lie after that target block. The
greedy choice takes the first eligible such block. The entire source
$B_D$ embeds in its available tail by Lemma~\ref{lem:blocks}; allowing
part of the source block to occur earlier cannot improve the endpoint
beyond this first eligible block. The remaining target suffix after the
greedy choice contains the suffix left by any later choice.

Inducting on the number of remaining source tokens proves completeness:
whenever a completion exists, the greedy step preserves the existence of
a completion. Conversely, each accepted finite step can be realized by
an actual letter match in a finite token or a periodic tail, and each
accepted infinite step by a cofinal embedding into the selected target
block. Concatenating these realizations gives an ordinal embedding, so
acceptance is sound.
\end{proof}

The target pointer never retreats. With bit operations and relation lookup
counted as unit cost, the decision procedure uses
$O(|u|_{\mathrm{tok}}+|v|_{\mathrm{tok}})$ time and constant auxiliary space.
For unbounded alphabet size, bitmask operations have their ordinary bit
costs; validating masks and the poset incurs additional cost. The
reference checker uses memoized exploration of every eligible target
endpoint rather than the first endpoint.

\subsection{Canonical quotient representatives}
\label{sec:canonical}

Proposition~\ref{prop:representation} produces finite expressions but, without
its normalization clause, not unique ones. The normalization is easy to carry
out on token expressions, and it makes the representatives canonical.

\paragraph{The normalization procedure.}
Process the tokens from left to right onto a stack. A letter token is pushed.
When an ideal token $[D]$ arrives, first pop the immediately preceding letter
tokens for as long as their letters belong to $D$, and then push $[D]$.
The final finite tail, that is the letters following the last ideal token, is
never altered. Equivalently: in each finite prefix preceding an ideal token
$[D]$, delete the longest suffix all of whose letters lie in $D$.

Each deleted suffix is absorbed by the ideal token that follows it, by
\eqref{eq:blockabsorb}, so the output is equivalent to the input; and the
output is exactly an expression of the shape \eqref{eq:normal}, since every
prefix is now empty or ends outside the downset of the block that follows it.
The procedure is linear in the number of tokens and is implemented as
\texttt{normalize} in \texttt{code/omega2.py}.

\begin{proposition}[Uniqueness of normalized expressions]\label{prop:canonical}
Two normalized token expressions over a finite poset $P$ are mutually
embeddable if and only if they are identical.
\end{proposition}
\begin{proof}
One direction is trivial. For the other, suppose $u$ and $v$ are normalized
and mutually embeddable. Ordinal length is monotone under embeddability, so
$u$ and $v$ have equal lengths; hence they have the same number $r$ of
canonical $\omega$-blocks and finite tails of the same length.

By Lemma~\ref{lem:equalblocks}, applied in both directions, the $i$th block of
$u$ and the $i$th block of $v$ are mutually embeddable, and so are the two
finite tails. Mutual embeddability of blocks forces mutual inclusion of their
downsets by Lemma~\ref{lem:blocks}, so the $i$th ideal tokens agree, say both
equal to $[D_i]$. Since both expressions are normalized, the finite prefixes
$p_i$ and $q_i$ preceding these tokens are empty or end outside $D_i$, so
\eqref{eq:fibre} turns the mutual embeddability of the blocks into mutual
embeddability of $p_i$ and $q_i$.

Finally, two finite words over a poset that embed into each other have equal
finite length, so both index maps are the identity; antisymmetry of $\le_P$
then forces equality of corresponding labels. Hence $p_i=q_i$ for every $i$,
the tails agree, and $u=v$.
\end{proof}

Consequently the quotient of $\Wordspace P$ has a countable, effectively
comparable system of canonical finite representatives, even though there are
uncountably many literal infinite words over an alphabet with at least two
elements. The forward reference promised after
Proposition~\ref{prop:representation} is discharged here.

\subsection{Why finite expansion is not an adequate substitute}
Replacing each $B_D$ by $M$ repetitions for a large integer $M$ changes the
problem. The absorption equivalence $aB_{\down a}\eqv B_{\down a}$ already
fails for the direction from length $M+1$ to length $M$ after a uniform
finite expansion. Conversely, finite approximations do not directly
certify the cofinal-boundary restrictions needed for reflection of product
maps. The supplied code therefore treats $\omega$-blocks symbolically and
uses the proven cofinal semantics exactly.

\section{Computational checks and reproducibility}
\label{sec:tests}

The included verification run passed \TestAssertions\ assertions. Its
random seed is \texttt{20260919}; all random checks supplement deterministic
checks. The report is stored in \texttt{data/verification.json}. The
recorded run used Python~3.14.4, and the code targets Python~3.10 or newer
without third-party packages. The suite is a single program: the counts below
come from one run, and they are not the sum of two separately reported runs.

\begin{table}[htbp]
\centering
\begin{tabular}{@{}lr@{}}
\toprule
Check & Number of assertions\\
\midrule
Exhaustive greedy/reference agreement, small posets & 38,107\\
Exhaustive greedy/reference agreement, binary domain & 726,242\\
Seeded longer-word greedy/reference agreement & 9,800\\
Normalization preserves every decided relation & 726,242\\
Normalization preserves the class and the ordinal length & 2,244\\
Mutually embeddable pairs have identical normal forms & 4,428\\
Exhaustive proper/universal product maps & 1,544,960\\
Binary support-family product maps & 44,667\\
Seeded longer-word product maps & 35,500\\
Padding reflection & 55,150\\
Padding has nonzero limit length & 4,024\\
Universal-only comparison formula & 20,640\\
Universal-only normalization & 2,660\\
Fixed-length product decomposition & 25,000\\
Maximal elements generate their downset & 356\\
Boundary, absorption, and regression cases & 16\\
Negative controls, required to fail & 4\\
\midrule
Total & \TestAssertions\\
\bottomrule
\end{tabular}
\caption{Executed symbolic checks. A ``word length'' in the test generator
means token length, not ordinal length.}
\label{tab:checks}
\end{table}

The exhaustive comparison of the two deciders uses all naturally labelled
posets of sizes zero through three and every source/target token word of
length at most two. The longer-word comparisons cover all naturally
labelled posets of sizes two through four and use token words of length
at most twelve.

\subsection{The exhaustive binary domain}
Over the two-letter alphabets the whole token domain up to length four is
enumerated and every ordered pair is tested. The antichain has five tokens,
giving $1+5+5^2+5^3+5^4=781$ expressions and $781^2=609{,}961$ ordered pairs;
the chain has four tokens, giving $341$ expressions and $116{,}281$ pairs, for
\BinaryPairs\ ordered pairs in all. For each pair the greedy and reference
comparators agree, normalization preserves the decided relation, and
mutually embeddable expressions have identical normal forms, which is the
executable form of Proposition~\ref{prop:canonical}.

Across the same two alphabets, every allowed support subfamily is used as an
insertion chain, giving \BinaryStages\ activation stages; the universal-only
exceptional activation is excluded, as the hypotheses require. Support
prefixes shared by different subfamilies are deliberately retested.

The three invalid shortcuts of Section~\ref{sec:negative} are executed as
negative controls: the suite asserts that the naive encodings do \emph{not}
reflect the componentwise order, so weakening a hypothesis would be detected
rather than silently accepted. The fixed-length decomposition of
Proposition~\ref{prop:exactproduct} is checked on randomly generated words of
prescribed length $\omega r+k$ over every naturally labelled poset of size at
most four.

For the product embeddings, every naturally labelled poset of sizes two
through four is processed, introducing all nonempty downsets in cardinality
order. There are \MarkerStages\ marker stages. At each stage the exhaustive
test uses a fixed four-word component pool and all pairs of tuples of
arities two and three. The seeded tests add tuples of arities two through
five with component token lengths at most eight. In total these give
\ProductChecks\ product-map comparisons. Separate padding tests allow all
one-token words, including already universal words, so padding reflection
is checked beyond the narrower family needed in the main induction.

The two embedding implementations share the mathematical periodic-tail
semantics. Agreement between them is a check of the greedy optimization,
not an independent derivation of that semantics. Likewise, no collection
of the finite assertions in Table~\ref{tab:checks} proves
$\omega^{\omega^4}$ to be the maximal order type. That conclusion rests on
the unbounded family of product embeddings and the ordinal calculation.

\subsection{Enumeration conventions}
The auxiliary enumerator considers relations on $\{0,\ldots,n-1\}$ whose
strict comparisons respect the usual numerical order. Every finite-poset
isomorphism class has such a labelling, but the enumeration is neither an
enumeration of all labellings nor an isomorphism-free enumeration. The
numbers of generated relations for $n=0,1,2,3,4,5$ are respectively
\[
 1,\ 1,\ 2,\ 7,\ 40,\ 357.
\]
These counts were computed by the included program. Histograms by the
number of downsets appear in \texttt{data/ideal\_counts.csv}. The $50$
individual posets of sizes one through four, with their predecessor and
successor masks, downset counts, and the exponent predicted by
\eqref{eq:main}, are listed in \texttt{data/finite\_posets.csv}; the row for
$n=1$ carries the literal entry \texttt{exception: omega\^{}2} rather than a
numerical exponent. Table examples are also exported to
\texttt{data/examples.csv}. No claim about uniform sampling of isomorphism
classes is intended.

To compute $j(P)$ for one small input poset, the code simply checks all
$2^n$ subsets for downward closure. This exponential enumeration is
sufficient for the examples and the proof tests. The article does not
claim a polynomial-time algorithm for downset counting.

\subsection{Reproduce the package}
From the root of the unpacked directory:
\begin{lstlisting}
python3 code/verify.py
latexmk -pdf -interaction=nonstopmode -halt-on-error article.tex
\end{lstlisting}
The Python command regenerates the verification reports and CSV tables.
It accepts \texttt{--seed}, \texttt{--samples}, \texttt{--pair-bound}, and
\texttt{--output-dir}; the defaults are $20260919$, $250$, $4$, and
\texttt{data/}. Raising \texttt{--pair-bound} enlarges the exhaustive binary
enumeration very quickly. A three-pass \texttt{pdflatex} build is also
supplied as \texttt{build.sh}.
The PDF can alternatively be built by running \texttt{pdflatex article.tex}
repeatedly until cross-references stabilize. No external bibliography
processor is required; a separate BibTeX file is supplied for reuse.

\section{Proof audit, scope, and further directions}
\label{sec:audit}

\subsection{Where a hidden gap would have to occur}
The proof is organized so its nonclassical burden is concentrated in a few
claims. The atom-count upper bound is a monotone-surjection argument.
Cofinal localization uses finiteness of the target token decomposition.
Synchronization uses the newness condition $D\nsubseteq E$ for all old
blocks, not merely distinctness of block labels. The proper-marker
reflection uses downward closure to exclude a guard from the next marker.
The full-marker reflection uses \emph{both} cancellation and the limit
length of the padded component.

These dependencies explain several tempting but incorrect shortcuts.
A source infinite block need not land wholly inside one target token.
The ordinary evaluation map from atom strings need not reflect order.
The universal marker admits no excluded finite letter. A finite test
cannot certify an infinite supremum. None of these shortcuts is used.

\paragraph{The four boundaries to check independently.}
Concretely, four points carry the argument, and a hidden gap would have to sit
at one of them.
\begin{enumerate}[label=\textnormal{(\arabic*)}]
\item Embeddings are \emph{strictly} increasing on ordinal positions.
      Without strictness, cofinal consumption of a target block fails and
      Lemma~\ref{lem:cofinal} is false as stated.
\item A newly introduced downset is contained in \emph{no} old downset,
      condition \eqref{eq:newness}. Mere distinctness is not enough, and
      \eqref{eq:negative} is the explicit counterexample.
\item Exclusion applies to the \emph{explicit pure periodic segment} of a
      displayed separator, not to the finite prefix that may precede it inside
      the same canonical block. This is Remark~\ref{rem:explicit}.
\item Full support requires an \emph{earlier proper} support, so that a guard
      of positive limit length can be built. Without one, the exceptional
      value \eqref{eq:exception} is the truth.
\end{enumerate}
The negative controls of Section~\ref{sec:negative} exhibit the failure of the
second and fourth safeguards, and of guarding altogether; the singleton
alphabet illustrates the resulting exceptional maximal order type.

\paragraph{A reader's guide.}
An independent referee can focus first on Theorems~\ref{thm:proper} and
\ref{thm:full}, or equivalently on their unified form
Theorem~\ref{thm:product}. Once their product-reflection assertions are
accepted, the classification follows by a short finite induction and the
classical finite-product formula, or by the elementary
Corollary~\ref{cor:altlevel} if one prefers to avoid that import. The binary
specialization in Section~\ref{sec:binary} provides a second presentation of
the same core argument, not an independent proof with different foundations.

Section~\ref{sec:exact} may be omitted entirely by a reader checking only the
resolution of the original two-letter question. It is the one part of the
article that uses the natural-product and natural-sum identities
Theorem~\ref{thm:classical}(iii) and~(iv) essentially, and nothing in
Sections~\ref{sec:intro}--\ref{sec:consequences} depends on it. Its role is to
document, in \eqref{eq:stratumsup}, why the answer cannot be obtained one
length at a time.

\subsection{Established within the argument versus not claimed}
Within this manuscript, the main theorem has a complete proof relative to
the classical results stated in Theorem~\ref{thm:classical}. There is no
remaining conjectural lemma inside that derivation. Nevertheless the
argument has not undergone independent expert review or proof-assistant
verification. The appropriate external status is therefore a
\emph{proposed solution}, with explicit hypotheses and reproducible
supporting checks, rather than a certified resolution.

The article does not compute all invariants of these quasi-orders: their
height, width, and exact isomorphism classification are different tasks.
It does not show that two alphabets with equal $n+j(P)$ produce isomorphic
word orders. It does not give a canonical representation for arbitrary
infinite input streams. It also does not assert that selected block
families at higher ordinal cutoffs have the same counting law.

\subsection{Why the proof stops at this cutoff}
The decisive compactness feature is that a target word below $\omega^2$
has only finitely many $\omega$-blocks. The image of a source $\omega$-block
must eventually stabilize in one of them. At length $\omega^2$, a source
$\omega$-block may instead visit infinitely many successive target
$\omega$-blocks, taking only finitely many points from each. Thus the
cofinal localization lemma in its present form is no longer valid when
that target length is allowed.

A higher-cutoff extension would need a hierarchy of separator ranks and a
new localization theorem that accounts for escape through infinitely many
lower-rank blocks. Counting indecomposable blocks without that theorem
would reproduce the same logical mistake that naive atom counting makes
at the universal-only stage.

For infinite alphabets, even when every individual word has finite image,
the set of possible recurrent downsets need not be finite, and a single
finite inclusion ordering does not exhaust the construction. The present
finite induction therefore does not directly transfer. These observations
identify precise missing mechanisms for extensions; they are not reasons
to reject further investigation.

\subsection{A concrete formal-verification target}
A useful formalization would first define ordinal concatenation and its
embedding relation for finite periodic-token codes, then prove the exact
symbolic decision theorem. The next targets are the cofinal boundary lemma,
proper-marker reflection, and universal-marker reflection. With the
classical maximal-order-type product and finite-word theorems available,
the final induction has very little additional complexity. The package
contains executable tests but no Lean development, and it makes no claim
of machine-checked proof.

\section{Conclusion}
\label{sec:conclusion}

The finite generator upper bound for words below $\omega^2$ is sharp for every
finite poset alphabet with at least two elements. The quantity that governs the
answer is the number of letters plus the number of nonempty downsets, and not
the alphabet's cardinality alone: the two-letter chain and the two-letter
antichain already separate.

The mechanism behind the sharpness is explicit. Each proper nonempty downset
can be activated as a separator by a terminal-letter guard, because a letter
outside a downward-closed set cannot be absorbed by the block that repeats it.
The full downset admits no such letter and can be activated only after an
earlier proper downset has supplied a guard of positive limit length, which no
finite initial segment of a periodic block can accommodate. When no earlier
proper downset is available the activation genuinely fails, and the
universal-only family has maximal order type $\Hn n\cdot\omega$ rather than
$\Hn{n+1}$; the singleton alphabet is the same phenomenon.

The proposed solution is therefore not merely a numerical lower bound matching
a known upper bound. It identifies and isolates the structural mechanism
responsible for the sharpness, in a form that is stated for arbitrary finite
posets, classified for arbitrary allowed families of periodic tails, made
effective by a canonical normal form and a linear-time symbolic comparator,
and audited both by worked negative controls and by a reproducible test suite.
Its external status remains that of an unrefereed proposed solution, in the
sense set out in Section~\ref{sec:audit}.

\appendix
\clearpage
\section{A compact self-contained statement of the construction}
\label{app:quickref}

This appendix restates the lower-bound mechanism on one page, in the
all-guarded form of Theorem~\ref{thm:product} and
Remark~\ref{rem:lastguard}. Fix a nonempty finite poset $P$, an allowed family
$\G\subseteq\J(P)\setminus\{\varnothing\}$, and a new nonempty downset
$D\notin\G$ with $D\nsubseteq E$ for every $E\in\G$. Set
\[
 g_D(u)=
 \begin{cases}
  uc,&D\ne P,\quad c\in P\setminus D,\\
  ucB_R,&D=P,\quad R\in\G\text{ proper},\quad c\in P\setminus R.
 \end{cases}
\]
The second line is defined only when the required $R$ exists; that proviso is
exactly the exceptional case of Theorem~\ref{thm:family}. For every
$m<\omega$ put
\[
 \Phi_m(\mathbf u)=g_D(u_0)B_Dg_D(u_1)B_D\cdots B_Dg_D(u_m).
\]
Then
\[
 \boxed{\quad
 \Phi_m(\mathbf u)\emb\Phi_m(\mathbf v)
 \quad\Longleftrightarrow\quad
 (\forall i\le m)\ u_i\emb v_i.
 \quad}
\]
The proof is synchronization (Lemma~\ref{lem:sync} with
Corollary~\ref{cor:between}), plus nonleakage (Lemma~\ref{lem:noleak}), plus
cancellation (Lemmas~\ref{lem:cancel} and~\ref{lem:padding}). If
$\mot(\W(P,\G))=\Hn q$ with $q=n+|\G|$, the equivalence yields
\[
 \Hn{q+1}=\sup_{m<\omega}\Hn q^{\,m+1}
 \ \le\ \mot\bigl(\W(P,\G\cup\{D\})\bigr)\ \le\ \Hn{q+1},
\]
where the left equality may be read either through the natural powers of
\eqref{eq:power}--\eqref{eq:sup} or through the ordinary powers of
\eqref{eq:amplification} and Lemma~\ref{lem:power}, and the right inequality
is Proposition~\ref{prop:upper}. This is the complete lower-bound mechanism
used in the article.

\clearpage
\section{A compact proof-dependency map}
\label{app:dependencies}

In outline, the argument has five layers:
\[
 \begin{gathered}
 \text{recurrent-downset representatives and the token upper bound}\\
 \Downarrow\\
 \text{cofinal synchronization and the two guards}\\
 \Downarrow\\
 \text{finite Cartesian-product embeddings}\\
 \Downarrow\\
 \left\{
 \begin{array}{c}
   \text{natural-product amplification (Thm.~\ref{thm:classical}(iii),
     \eqref{eq:power}--\eqref{eq:sup})}\\
   \text{or elementary ordinary-power amplification
     (Lem.~\ref{lem:power}, \eqref{eq:amplification})}
 \end{array}
 \right\}\\
 \Downarrow\\
 \text{support-family classification and the finite-poset formula.}
 \end{gathered}
\]
The fourth layer is the only one at which the article offers two independent
routes; either suffices for Theorem~\ref{thm:main}. The table below expands
the same dependencies statement by statement.

\begin{longtable}{@{}>{\raggedright\arraybackslash}p{0.29\textwidth}>{\raggedright\arraybackslash}p{0.31\textwidth}>{\raggedright\arraybackslash}p{0.32\textwidth}@{}}
\toprule
Statement & Inputs & Role\\
\midrule
\endhead
Finite representation, Proposition~\ref{prop:representation}
 & Finiteness of $P$; recurrent labels; monotone concatenation
 & Identifies the full word quotient with a finite-token quotient.\\[3pt]
Atom-count upper bound, Proposition~\ref{prop:upper}
 & Finite-word formula; monotone surjection
 & Gives $\mot(\W(P,\F))\le H_{n+|\F|}$.\\[3pt]
Cofinal localization, Lemma~\ref{lem:cofinal}
 & Finitely many target tokens; strict increase; periodic recurrence
 & Finds the target block used cofinally by each source marker.\\[3pt]
Synchronization, Lemma~\ref{lem:sync}
 & Cofinal localization; no old block contains the new downset
 & Forces the $i$th new marker to correspond to the $i$th new marker.\\[3pt]
Proper product embedding, Theorem~\ref{thm:proper}
 & Synchronization; a letter outside the proper downset
 & Embeds every finite Cartesian power of the old stage.\\[3pt]
Limit padding, Lemma~\ref{lem:padding}
 & A proper old block; an excluded letter; final-letter cancellation
 & Produces an order-reflecting map into nonzero limit lengths.\\[3pt]
Universal product embedding, Theorem~\ref{thm:full}
 & Synchronization; limit padding; the finite-bound contradiction
 & Adds the universal block without losing the product lower bound.\\[3pt]
Family classification, Theorem~\ref{thm:family}
 & The two product maps; classical natural product; atom upper bound
 & Computes every selected family, including its sole exception.\\[3pt]
Main theorem, Theorem~\ref{thm:main}
 & Full family has a proper downset when $n\ge2$; finite representation
 & Gives $\omega^{\omega^{n+j(P)-2}}$ and the two boundary cases.\\[3pt]
Unified product embedding, Theorem~\ref{thm:product}
 & The same synchronization, with one case-split guard map
 & Restates Theorems~\ref{thm:proper} and~\ref{thm:full} as one equivalence.\\[3pt]
Elementary amplification, Lemma~\ref{lem:power} and
Corollary~\ref{cor:altlevel}
 & A maximal linear extension; the last-unequal-coordinate order
 & Replaces the natural-product import in the main line.\\[3pt]
One-block normal form, Lemma~\ref{lem:oneblock}
 & Last occurrence outside the recurrent downset
 & Supplies the admissible prefixes used by \eqref{eq:fibre}.\\[3pt]
Equal-count rigidity, Lemma~\ref{lem:equalblocks}
 & Cofinal localization at a fixed block count
 & Whole-block containment; used only in
   Sections~\ref{sec:exact} and~\ref{sec:canonical}.\\[3pt]
Exact-length strata, Propositions~\ref{prop:exactomega}
and~\ref{prop:exactproduct}
 & Equal-count rigidity; natural sum and product
 & Supplementary; yields the warning \eqref{eq:stratumsup}.\\[3pt]
Canonical representatives, Proposition~\ref{prop:canonical}
 & Normalized prefixes; equal-count rigidity; \eqref{eq:fibre}
 & Makes the finite representation unique and effective.\\
\bottomrule
\end{longtable}

\section{Minimal worked use of the code}
\label{app:code}

The following example can be run with \texttt{code/} on the Python import
path. Downset masks \texttt{1}, \texttt{2}, and \texttt{3} represent
$\{a\}$, $\{b\}$, and $\{a,b\}$.

\begin{lstlisting}[language=Python]
from omega2 import (FinitePoset, letter, omega, embeds, normalize,
                    encode_product, theorem_value)

P = FinitePoset.antichain(2)
a, b = letter(0), letter(1)
A, B, U = omega(1), omega(2), omega(3)

assert embeds(P, (a, U), (U,))       # finite-prefix absorption
assert embeds(P, (A,), (U,))        # pure tail into universal tail
assert not embeds(P, (A, B), (U,))  # two cofinal blocks cannot share one
assert not embeds(P, (U,), (A, B))  # no eligible cofinal destination

assert normalize((a, U)) == (U,)    # canonical quotient representative
assert normalize((b, A)) == (b, A)  # b is outside the downset of A

# The unified guarded product map, introducing B over the family {A}.
print(encode_product(P, (1,), 2, ((), (a,))))
# (('f', 0), ('w', 2), ('f', 0), ('f', 0))

print(theorem_value(P))
# omega^(omega^4)

C = FinitePoset.chain(2)
print(theorem_value(C))
# omega^(omega^3)

print(theorem_value(P, family=(P.full,)))
# omega^(omega^1 + 1) = omega^(omega + 1)
\end{lstlisting}

The last output illustrates the exceptional universal-only family. Its
maximal order type is much smaller than $H_3=\omega^{\omega^2}$, even
though the two finite letters and the universal block form three distinct
formal tokens.

\section{Literature and reproducibility record}
The chosen source was checked both in arXiv HTML and in the PDF rendering;
Example~3.15 is on printed page~10, which is zero-based PDF page~9. Its
version identifier is \texttt{2409.03199v2}. The source of the open question,
the two standard maximal-order-type references, and the distinction between
proven classical inputs and the present proposed argument are recorded in
\texttt{notes/literature\_audit.md}.

The archive contains this source and its compiled PDF, a reusable
\texttt{references.bib}, the symbolic module, the reproducible verification
script, machine-readable test results, example and downset-count tables,
a proof-audit note, and a README. It contains no third-party article PDFs,
font files, checksum files, or LaTeX build intermediates. The file inventory
is the following.

\begin{longtable}{@{}p{0.38\linewidth}p{0.56\linewidth}@{}}
\toprule
File & Purpose\\
\midrule
\endhead
\texttt{article.tex} & Complete editable source of this manuscript.\\
\texttt{article.pdf} & Compiled article.\\
\texttt{references.bib} & Bibliographic records, maintained alongside the
  embedded \texttt{thebibliography}.\\
\texttt{build.sh} & Three-pass \texttt{pdflatex} build script.\\
\texttt{README.md} & Scope, result list, reproduction, and library usage.\\
\texttt{code/omega2.py} & Finite posets, downsets, exact token comparators,
  normalization, the guard maps, and the product encodings.\\
\texttt{code/verify.py} & Exhaustive and seeded checks, negative controls,
  and table generation.\\
\texttt{data/verification.json}, \texttt{data/verification.txt}
  & Recorded output of the default verification run.\\
\texttt{data/test\_summary.tex} & Test-count macros used by this article.\\
\texttt{data/examples.csv} & The computed finite-poset examples of
  Table~\ref{tab:examples}.\\
\texttt{data/ideal\_counts.csv} & Downset-count histogram for naturally
  labelled posets of size at most five.\\
\texttt{data/finite\_posets.csv} & The $50$ naturally labelled posets through
  size four, individually, with predecessor and successor masks, downset
  counts, and the exponent predicted by \eqref{eq:main}.\\
\texttt{notes/proof\_audit.md} & Hypotheses, critical steps, and known false
  shortcuts.\\
\texttt{notes/literature\_audit.md} & Source, version, and literature-search
  scope.\\
\bottomrule
\end{longtable}

\begin{thebibliography}{9}
\bibitem{Altman2026}
Harry Altman.
\newblock \emph{Bounding finite-image sequences of length $\omega^k$}.
\newblock arXiv:2409.03199v2, revised March 10, 2026; manuscript dated
March 9, 2026. In particular, Example 3.15, p.~10.
\newblock \url{https://arxiv.org/abs/2409.03199v2}.

\bibitem{deJonghParikh1977}
D.~H.~J. de Jongh and Rohit Parikh.
\newblock Well-partial orderings and hierarchies.
\newblock \emph{Indagationes Mathematicae (Proceedings)}, 80(3):195--207,
1977; also cited as \emph{Indagationes Mathematicae} 39.
\newblock \url{https://doi.org/10.1016/1385-7258(77)90067-1}.

\bibitem{Schmidt2020}
Diana Schmidt.
\newblock Well-Partial Orderings and their Maximal Order Types.
\newblock Habilitationsschrift, Heidelberg, 1979.
\newblock Reprinted in P.~Schuster, M.~Seisenberger, and A.~Weiermann, editors,
\emph{Well-Quasi Orders in Computation, Logic, Language and Reasoning},
Trends in Logic 53, pp.~351--391. Springer, 2020.
\newblock \url{https://doi.org/10.1007/978-3-030-30229-0_13}.

\bibitem{Higman1952}
Graham Higman.
\newblock Ordering by divisibility in abstract algebras.
\newblock \emph{Proceedings of the London Mathematical Society}, series~3,
\textbf{2}(1):326--336, 1952.
\newblock \url{https://doi.org/10.1112/plms/s3-2.1.326}.

\bibitem{ChopraPakhomov2026}
Alakh Dhruv Chopra and Fedor Pakhomov.
\newblock \emph{Well-quasi-orders on finite trees and transfinite sequences}.
\newblock arXiv:2602.09830v1, February 10, 2026.
\newblock \url{https://arxiv.org/abs/2602.09830v1}.
\end{thebibliography}

\end{document}
